\usepackage[utf8]{inputenc}
\usepackage[T1]{fontenc}
\usepackage[portuguese]{babel}

% --- FONTE PADRÃO (Estilo Didático Encorpado) ---
\usepackage{helvet} 
\renewcommand{\familydefault}{\sfdefault}

% --- GEOMETRIA E ESPAÇAMENTO ---
\usepackage[top=1.6cm, bottom=1.5cm, left=0.9cm, right=0.9cm, headheight=22pt, headsep=0.4cm]{geometry}
\usepackage{microtype} 
\usepackage{setspace}
\setstretch{1.0}
\usepackage{parskip}
\setlength{\parskip}{2pt}
\setlength{\parindent}{0pt}

\newlength{\espacoPosTitulo}\setlength{\espacoPosTitulo}{0.3cm}
\newlength{\espacoPreQuestao}\setlength{\espacoPreQuestao}{0.1cm}
\newlength{\espacoQuestao}\setlength{\espacoQuestao}{0.2cm}
\newlength{\espacoAlt}\setlength{\espacoAlt}{0.2cm}

\usepackage{enumitem}
\setlist{itemsep=0.4cm, topsep=0.1cm, parsep=0pt, partopsep=0pt}
\newcommand{\larguraTikz}{0.85\linewidth} 
% --- TAG SEMÂNTICA PARA LEITURA DE IA (INVISÍVEL NO PDF) ---
\newenvironment{enunciadoLiteral}{}{}

% --- MATEMÁTICA E CORES ---
\usepackage{amsmath, amsfonts, amssymb, nccmath, mathastext}
\usepackage{xcolor}
\definecolor{indigo}{HTML}{4F46E5}
\definecolor{CorTitUm}{HTML}{FF6F61}
\definecolor{CorTitDois}{HTML}{FF6F61}
\definecolor{CorTitTres}{HTML}{658894}
\definecolor{CorLinha}{HTML}{B0B0B0} 
\definecolor{metal}{HTML}{7F8C8D}
\definecolor{vidro}{HTML}{3498DB}
\definecolor{ouro}{HTML}{F1C40F}
\definecolor{concreto}{HTML}{95A5A6}
\definecolor{madeira}{HTML}{D35400}
\definecolor{CinzaEscuro}{HTML}{4A4A4A} 
\definecolor{CinzaMedio}{HTML}{848484} 
\definecolor{Cinzablack}{HTML}{353535} 

% --- MINI-CABEÇALHO E RODAPÉ AUTORAL (TWOSIDE) ---
\usepackage{fancyhdr}
\pagestyle{fancy}
\fancyhf{} 
\renewcommand{\headrulewidth}{0.3pt} 
\renewcommand{\footrulewidth}{0.3pt}
\fancyhead[LE]{\small\sffamily\bfseries\color{gray!75} EQUAÇÕES E INEQUAÇÕES DO 1º GRAU}
\ifgabarito
    \fancyhead[RO]{\small\sffamily\bfseries\color{CorTitUm}{LIVRO DO PROFESSOR -- SUPLEMENTAR}}
\else
    \fancyhead[RO]{\small\sffamily\color{gray!75} \texttt{www.profraphaelpaulino.com.br}}
\fi
\fancyfoot[LE]{\small \textbf{\thepage}}
\fancyfoot[RO]{\small \textbf{\thepage}}

% --- CAIXAS DE DESTAQUE ---
\usepackage[most]{tcolorbox}
\newtcolorbox{boxresponda}{colback=white, colframe=gray!45, boxrule=0pt, leftrule=1.7pt, sharp corners, enlarge left by=0.4cm, width=\linewidth-0.4cm, left=8pt, right=0pt, top=2pt, bottom=2pt, fontupper=\small}
\definecolor{CorDicaBorda}{HTML}{17c1be} 
\definecolor{CorDicaFundo}{HTML}{f3fbfb} 
\newtcolorbox{boxdica}{colback=CorDicaFundo, colframe=CorDicaBorda, boxrule=0pt, leftrule=2.5pt, sharp corners, width=\linewidth, left=8pt, right=8pt, top=6pt, bottom=6pt, halign=center, fontupper=\small}

% --- NOVA CAIXA PEDAGÓGICA E CORES (NOTA VERDE) ---
\definecolor{CorNotaBorda}{HTML}{10B981} 
\definecolor{CorNotaFundo}{HTML}{F0FDF4} 
\newtcolorbox{boxexplicacao}{colback=CorNotaFundo, colframe=CorNotaBorda, boxrule=0pt, leftrule=2.5pt, sharp corners, width=\linewidth, left=8pt, right=8pt, top=6pt, bottom=6pt, fontupper=\small}

% --- TÍTULOS ---
\usepackage{titlesec}
\titleformat{\section}{\color{black}\large\bfseries}{\thesection}{0em}{\vspace{0.1cm}\titlerule[0.8pt]}
\titlespacing*{\section}{0pt}{10pt}{6pt} 
\titleformat{\subsubsection}[runin]{\color{black}\bfseries}{}{0em}{}
\titlespacing*{\subsubsection}{0pt}{0pt}{0.15cm}

% --- COLUNAS E GRÁFICOS ---
\usepackage{multicol}
\setlength{\columnsep}{1cm}
\setlength{\columnseprule}{0.5pt}
\def\columnseprulecolor{\color{CorLinha}}
\usepackage{adjustbox, tikz, pgfplots, circuitikz}
\usetikzlibrary{shadows, shapes.misc, positioning, calc}
\pgfplotsset{compat=1.18}

\begin{document}
\thispagestyle{empty}
\color{Cinzablack}
\vspace*{-1.8cm}

% --- CABEÇALHO PRINCIPAL AUTORAL (Apenas 1ª página) ---
\noindent
\begin{minipage}{\textwidth}
    \fontfamily{phv}\selectfont
    \noindent
    \begin{minipage}[t]{0.65\linewidth}
        {\Large \bfseries \MakeUppercase{Caderno de Atividades Suplementares}}\\[0.1cm]
        {\large \textmd{\textcolor{CinzaEscuro}{Unidade: EQUAÇÕES E INEQUAÇÕES DO 1º GRAU}}}
    \end{minipage}%
    \begin{minipage}[t]{0.35\linewidth}
        \raggedleft
        \ifgabarito
            {\large \bfseries \textcolor{CorTitUm}{[PROFESSOR]}}\\[0.05cm]
        \fi
        \small \textbf{Prof. Raphael Paulino}\\
        \textsf{\small \textcolor{indigo}{\texttt{profraphaelpaulino.com.br}}}
    \end{minipage}
    
    \vspace{0.15cm}
    \noindent\rule{\linewidth}{1.2pt}
    \vspace{-0.1cm}
    \footnotesize\textsf{\textbf{ÁREA:} MATEMÁTICA \hfill \textbf{NÍVEL:} 7º ANO}
    \vspace{0.1cm}
    \noindent\rule{\linewidth}{0.4pt}
\end{minipage}

\par\vspace{0.3cm} 
\raggedcolumns
\begin{multicols*}{2}
\vspace{0.1cm}

\subsection*{Capítulo 1 e 2: Conceitos e Resolução de Equações}
\vspace{-0.2cm}\noindent\rule{\linewidth}{1.5pt} 
\par\vspace{\espacoPosTitulo}
\subsubsection*{1.} 
\begin{enunciadoLiteral}
A figura ilustra uma balança de pratos clássica em perfeito estado de equilíbrio. No prato da esquerda, foram colocados três blocos metálicos de massa \textbf{x} e um peso padrão aferido em \textbf{500 g}. No prato da direita, encontra-se um único peso padrão de \textbf{2 kg}.

\begin{center}
% ==========================================
% CONTROLE INDIVIDUAL DESTA IMAGEM:
% 1. CORTAR BORDAS: trim=1cm 0cm 1cm 0cm
% 2. TAMANHO: \larguraTikz por um valor
% ==========================================
\begin{adjustbox}{trim=0cm 0cm 0cm 0cm, clip, max width=\larguraTikz}
\begin{tikzpicture}
% --- APLICANDO DROP SHADOW E CORES ---
\definecolor{baseScale}{HTML}{2C3E50}
\definecolor{prato}{HTML}{95A5A6}
\definecolor{pesoFixo}{HTML}{C0392B}
\definecolor{pesoIncognito}{HTML}{F39C12}

% LAYER 1: Base e Eixo Central
\fill[baseScale, drop shadow={opacity=0.2, shadow xshift=1pt, shadow yshift=-2pt}, rounded corners=2pt] (-0.5,0) rectangle (0.5,0.4);
\fill[baseScale, drop shadow={opacity=0.2, shadow xshift=1pt, shadow yshift=-2pt}] (-0.1,0.4) rectangle (0.1,3);

% LAYER 2: Braço de equilíbrio
\fill[baseScale!80, drop shadow={opacity=0.2, shadow xshift=1pt, shadow yshift=-2pt}, rounded corners=1pt] (-3,2.8) rectangle (3,3.0);
\fill[baseScale] (0,2.9) circle (0.15);

% LAYER 3: Fios e Pratos
\draw[thick, gray] (-2.9,2.8) -- (-3.5,1.5);
\draw[thick, gray] (-2.9,2.8) -- (-2.3,1.5);
\draw[thick, gray] (2.9,2.8) -- (2.3,1.5);
\draw[thick, gray] (2.9,2.8) -- (3.5,1.5);

\fill[prato, drop shadow={opacity=0.2, shadow xshift=1pt, shadow yshift=-2pt}] (-3.7,1.4) to[out=-20,in=200] (-2.1,1.4) -- (-2.1,1.5) -- (-3.7,1.5) -- cycle;
\fill[prato, drop shadow={opacity=0.2, shadow xshift=1pt, shadow yshift=-2pt}] (2.1,1.4) to[out=-20,in=200] (3.7,1.4) -- (3.7,1.5) -- (2.1,1.5) -- cycle;

% LAYER 4: Pesos Prato Esquerdo (3x + 500g)
\fill[pesoIncognito, drop shadow={opacity=0.3, shadow xshift=1pt, shadow yshift=-1pt}, rounded corners=1pt] (-3.4,1.5) rectangle (-2.9,2.0);
\node[white, font=\bfseries\small] at (-3.15,1.75) {x};
\fill[pesoIncognito, drop shadow={opacity=0.3, shadow xshift=1pt, shadow yshift=-1pt}, rounded corners=1pt] (-2.8,1.5) rectangle (-2.3,2.0);
\node[white, font=\bfseries\small] at (-2.55,1.75) {x};
\fill[pesoIncognito, drop shadow={opacity=0.3, shadow xshift=1pt, shadow yshift=-1pt}, rounded corners=1pt] (-3.1,2.0) rectangle (-2.6,2.5);
\node[white, font=\bfseries\small] at (-2.85,2.25) {x};

\fill[pesoFixo, drop shadow={opacity=0.3, shadow xshift=1pt, shadow yshift=-1pt}] (-2.2,1.5) -- (-2.0,1.8) -- (-1.8,1.5) -- cycle;
\fill[pesoFixo] (-2.1,1.8) rectangle (-1.9,1.9);
\node[fill=white, inner sep=1pt, font=\tiny\bfseries, rounded corners=1pt] at (-2.0,2.1) {500g};

% LAYER 5: Peso Prato Direito (2kg)
\fill[pesoFixo, drop shadow={opacity=0.3, shadow xshift=2pt, shadow yshift=-1pt}] (2.5,1.5) -- (2.9,2.2) -- (3.3,1.5) -- cycle;
\fill[pesoFixo] (2.7,2.2) rectangle (3.1,2.4);
\node[fill=white, inner sep=1pt, font=\bfseries\small, rounded corners=1pt] at (2.9,2.7) {2 kg};

\end{tikzpicture}
\end{adjustbox}
\end{center}

Com base no princípio da equivalência, determine a massa matemática exata de um único bloco \textbf{x}, expressa em gramas. 
\end{enunciadoLiteral}

\ifgabarito
    % --- CAIXA LARANJA (Resolução e Tutor) ---
    \begin{tcolorbox}[colback=CorTitUm!5, colframe=CorTitUm, boxrule=1pt, arc=4pt, left=6pt, right=6pt, top=6pt, bottom=6pt]
    \small\textbf{\textcolor{CorTitUm}{Roteiro do Professor e Tutor IA:}}\par\vspace{0.1cm}
    \color{CinzaEscuro}
    \textbf{1. Ancoragem (Abordagem Socrática):} Antes de montar qualquer equação, pergunte ao aluno: "Se a balança está equilibrada, o peso total da esquerda é igual ou diferente do da direita? Como escrevemos 'equilíbrio' na matemática? Podemos trabalhar com gramas e quilogramas misturados na mesma conta?"\par\vspace{0.1cm}
    
    \textbf{2. Micro-passos da Resolução:}\\
    \textit{Passo 1:} Unificação de grandezas. Converter \textbf{2 kg} para gramas $\rightarrow$ $2 \text{ kg} = 2000 \text{ g}$.\\
    \textit{Passo 2:} Tradução da imagem para linguagem algébrica. O prato esquerdo tem três massas \textbf{x} mais \textbf{500 g}. O equilíbrio vira o sinal de igualdade $\rightarrow$ $3x + 500 = 2000$.\\
    \textit{Passo 3:} Isolar o termo desconhecido subtraindo 500 dos dois lados $\rightarrow$ $3x = 2000 - 500 \Rightarrow 3x = 1500$.\\
    \textit{Passo 4:} Dividir o resultado pela quantidade de blocos $\rightarrow$ $x = \mfrac{1500}{3} \Rightarrow x = 500 \text{ g}$.\par\vspace{0.1cm}
    
    \textbf{3. Mapeamento de Armadilhas (Alerta Tutor):} O erro mais clássico aqui é o aluno ignorar a conversão de unidades e montar a equação $3x + 500 = 2$. Isso geraria um valor negativo para a massa. Se o aluno travar, questione sobre quantos gramas cabem em um quilo.
    \end{tcolorbox}
    \vspace{0.15cm}
    
    % --- CAIXA VERDE (Porto Seguro Didático do Professor) ---
    \begin{boxexplicacao}
    \textbf{Porto Seguro Didático (O "Pulo do Gato"):} Uma equação do 1º grau nada mais é do que uma gangorra nivelada. Tudo o que você tira ou coloca de um lado, é OBRIGADO a tirar ou colocar do outro na mesma proporção para a gangorra não despencar. O sinal de igual ($=$) é o eixo central que sustenta o equilíbrio. E nunca se esqueça: não podemos somar "bananas com laranjas" — ou seja, converta as unidades para a mesma base antes de calcular!
    \end{boxexplicacao}
\fi
    
\par\vspace{\espacoQuestao}


\noindent\textcolor{CorLinha}{\rule{\linewidth}{0.5pt}} 
\par\vspace{\espacoPreQuestao}
\subsubsection*{2.} 
\begin{enunciadoLiteral}
Em laboratórios de engenharia de materiais, a temperatura de fusão ideal de uma liga de alumínio é calibrada por meio de uma constante \textbf{T}. O cálculo rigoroso para encontrar essa constante obedece à modelagem matemática:

$ 5(T - 12) + 3T = 8T - 4(15 - 2T) $

Execute as manipulações algébricas necessárias, aplicando a propriedade distributiva, e determine o valor da constante \textbf{T}.
\end{enunciadoLiteral}

\ifgabarito
    % --- CAIXA LARANJA (Resolução e Tutor) ---
    \begin{tcolorbox}[colback=CorTitUm!5, colframe=CorTitUm, boxrule=1pt, arc=4pt, left=6pt, right=6pt, top=6pt, bottom=6pt]
    \small\textbf{\textcolor{CorTitUm}{Roteiro do Professor e Tutor IA:}}\par\vspace{0.1cm}
    \color{CinzaEscuro}
    \textbf{1. Ancoragem (Abordagem Socrática):} Pergunte: "Temos parênteses nesta equação. Quem tem prioridade de resolução? Como fazemos para 'quebrar' os parênteses multiplicando o termo de fora por todos os de dentro?"\par\vspace{0.1cm}
    
    \textbf{2. Micro-passos da Resolução:}\\
    \textit{Passo 1:} Aplicar a propriedade distributiva no primeiro membro $\rightarrow$ $5T - 60 + 3T$.\\
    \textit{Passo 2:} Aplicar a distributiva no segundo membro (atenção redobrada ao sinal negativo multiplicando outro negativo) $\rightarrow$ $8T - 60 + 8T$.\\
    \textit{Passo 3:} Juntar os termos semelhantes em cada lado $\rightarrow$ $8T - 60 = 16T - 60$.\\
    \textit{Passo 4:} Organizar as variáveis de um lado e números do outro $\rightarrow$ $8T - 16T = -60 + 60$.\\
    \textit{Passo 5:} Solução final $\rightarrow$ $-8T = 0 \Rightarrow T = 0$.\par\vspace{0.1cm}
    
    \textbf{3. Mapeamento de Armadilhas (Alerta Tutor):} O erro principal está no segundo membro: ao fazer $-4 \cdot (-2T)$, muitos alunos colocam $-8T$ em vez de $+8T$. Se o aluno errar, guie-o a revisar a regra de sinais da multiplicação. Outro erro comum é chegar em $-8T = 0$ e concluir que "não tem solução", confundindo com divisão por zero.
    \end{tcolorbox}
    \vspace{0.15cm}
    
    % --- CAIXA VERDE (Porto Seguro Didático do Professor) ---
    \begin{boxexplicacao}
    \textbf{Porto Seguro Didático (O "Pulo do Gato"):} O famoso "chuveirinho" (propriedade distributiva) molha todo mundo que está dentro do box (os parênteses). Lembre o aluno de que o sinal que acompanha o número do lado de fora também entra na dança! E se chegar ao final e der zero de um lado da balança ($8T = 0$), fique calmo: o número zero é uma resposta válida e respeitável, só significa que a variável não tem peso.
    \end{boxexplicacao}
\fi

\par\vspace{\espacoQuestao}


\noindent\textcolor{CorLinha}{\rule{\linewidth}{0.5pt}}
\par\vspace{\espacoPreQuestao}
\subsubsection*{3.} 
\begin{enunciadoLiteral}
Um estudante de robótica precisava calcular a tensão \textbf{x} de um circuito usando a seguinte equação inicial:

$ 4x - (x + 6) = 2(x - 1) $

O estudante realizou os cálculos e apresentou o seguinte registro no seu caderno de anotações:

\textit{Linha 1:} $ 4x - x + 6 = 2x - 2 $ \\
\textit{Linha 2:} $ 3x + 6 = 2x - 2 $ \\
\textit{Linha 3:} $ 3x - 2x = -2 + 6 $ \\
\textit{Linha 4:} $ x = 4 $

Houve uma falha conceitual severa durante a execução dos procedimentos. Analise minuciosamente o caderno do estudante.

\begin{boxresponda}
\textbf{Responda:} Em qual linha o estudante cometeu o primeiro erro matemático? Justifique textualmente o que causou o equívoco e apresente a resolução correta a partir daquele ponto.
\end{boxresponda}
\end{enunciadoLiteral}

\ifgabarito
    % --- CAIXA LARANJA (Resolução e Tutor) ---
    \begin{tcolorbox}[colback=CorTitUm!5, colframe=CorTitUm, boxrule=1pt, arc=4pt, left=6pt, right=6pt, top=6pt, bottom=6pt]
    \small\textbf{\textcolor{CorTitUm}{Roteiro do Professor e Tutor IA:}}\par\vspace{0.1cm}
    \color{CinzaEscuro}
    \textbf{1. Ancoragem (Abordagem Socrática):} Peça ao aluno para não olhar os números grandes e focar apenas no sinal de "menos" antes do parênteses. "Quando há um sinal negativo na frente de uma chave ou parêntese, o que ele faz com quem mora lá dentro?"\par\vspace{0.1cm}
    
    \textbf{2. Micro-passos da Resolução:}\\
    \textit{Passo 1:} Diagnóstico do erro. O erro ocorreu logo na \textbf{Linha 1}. O estudante não distribuiu o sinal negativo para o $+6$ que estava dentro dos parênteses. Ele apenas retirou os parênteses, escrevendo $+6$ em vez de $-6$.\\
    \textit{Passo 2:} Justificativa. Um sinal negativo na frente do parênteses indica o "oposto" de tudo que está lá dentro, funcionando como se fosse uma multiplicação por $-1$.\\
    \textit{Passo 3:} Resolução corrigida (Linha 1 correta) $\rightarrow$ $4x - x - 6 = 2x - 2$.\\
    \textit{Passo 4:} Redução (Linha 2) $\rightarrow$ $3x - 6 = 2x - 2$.\\
    \textit{Passo 5:} Agrupamento (Linha 3) $\rightarrow$ $3x - 2x = -2 + 6$.\\
    \textit{Passo 6:} Solução Final $\rightarrow$ $x = 4$.\par\vspace{0.1cm}
    
    \textbf{3. Mapeamento de Armadilhas (Alerta Tutor):} Ironicamente, o erro de sinal na Linha 1 e o erro de agrupamento na Linha 3 (passar o $+6$ como $+6$ e não como $-6$) acabaram se anulando no resultado, mas o processo está duplamente errado. Se o aluno não perceber, peça para ele substituir $x = 4$ na equação original e provar se a balança equilibra (não equilibra).
    \end{tcolorbox}
    \vspace{0.15cm}
    
    % --- CAIXA VERDE (Nota Pedagógica) ---
    \begin{boxexplicacao}
    \textbf{Justificativa Curricular (Raio-X da Questão):} Avalia a capacidade metacognitiva do aluno (aprender com o erro do outro). Identificar falhas processuais sedimenta melhor a regra de distribuição de sinais do que resolver dezenas de contas mecanicamente.
    \end{boxexplicacao}

% --- CAIXA VERDE (Porto Seguro Didático do Professor) ---
    \begin{boxexplicacao}
    \textbf{Porto Seguro Didático (O "Pulo do Gato"):} O sinal de "menos" colado no parênteses age como um "inversor universal". Ele entra na casa (os parênteses) e inverte a polaridade elétrica de todos os moradores. Se era positivo, fica negativo. Se era negativo, vira positivo. Esquecer de inverter o último termo é a principal causa de acidentes nas equações.
    \end{boxexplicacao}
\fi

\par\vspace{\espacoQuestao}


\noindent\textcolor{CorLinha}{\rule{\linewidth}{0.5pt}} 
\par\vspace{\espacoPreQuestao}
\subsubsection*{4.} 
\begin{enunciadoLiteral}
Durante a execução de um projeto de urbanismo, a divisão de um terreno foi calculada de forma que a soma da metade da área total (\textbf{x}) com sua terça parte seja igual a exatos 1000 metros quadrados, gerando a seguinte modelagem:

$ \mfrac{x}{2} + \mfrac{x}{3} = 1000 $

Com base nos princípios de resolução de equações envolvendo coeficientes fracionários, avalie as alternativas a seguir e identifique o valor da área total \textbf{x}. Em seguida, justifique obrigatoriamente o erro conceitual presente na alternativa B.

\begin{enumerate}[label=\alph*), itemsep=\espacoAlt]
\item $ x = 500 $
\item $ x = 200 $
\item $ x = 1200 $
\item $ x = 6000 $
\end{enumerate}
\end{enunciadoLiteral}

\ifgabarito
    % --- CAIXA LARANJA (Resolução e Tutor) ---
    \begin{tcolorbox}[colback=CorTitUm!5, colframe=CorTitUm, boxrule=1pt, arc=4pt, left=6pt, right=6pt, top=6pt, bottom=6pt]
    \small\textbf{\textcolor{CorTitUm}{Roteiro do Professor e Tutor IA:}}\par\vspace{0.1cm}
    \color{CinzaEscuro}
    \textbf{1. Ancoragem (Abordagem Socrática):} "Você pode somar duas frações que têm fatias de tamanhos diferentes (denominadores 2 e 3)? Como encontramos um tamanho de fatia comum para toda a equação?"\par\vspace{0.1cm}
    
    \textbf{2. Micro-passos da Resolução:}\\
    \textit{Passo 1:} Encontrar o Mínimo Múltiplo Comum (MMC) dos denominadores 2, 3 e 1 (que está oculto sob o 1000). O MMC é 6.\\
    \textit{Passo 2:} Multiplicar TODOS os termos da equação pelo MMC para eliminar as frações $\rightarrow$ $6 \cdot (\mfrac{x}{2}) + 6 \cdot (\mfrac{x}{3}) = 6 \cdot (1000)$.\\
    \textit{Passo 3:} Simplificar as frações $\rightarrow$ $3x + 2x = 6000$.\\
    \textit{Passo 4:} Juntar os termos semelhantes $\rightarrow$ $5x = 6000$.\\
    \textit{Passo 5:} Isolar \textbf{x} $\rightarrow$ $x = \mfrac{6000}{5} \Rightarrow x = 1200$. A resposta correta é a letra \textbf{C}.\\
    \textit{Passo 6:} Refutação da Letra B. O erro da alternativa B ($x = 200$) ocorre quando o aluno soma os denominadores de forma errada ($2+3=5$) e monta $\mfrac{x}{5} = 1000$, esquecendo-se da regra básica de frações de tirar o MMC, e depois divide 1000 por 5.\par\vspace{0.1cm}
    
    \textbf{3. Mapeamento de Armadilhas (Alerta Tutor):} O erro mais clássico em equações fracionárias é o aluno tirar o MMC apenas do lado esquerdo (onde estão as letras) e esquecer de multiplicar o lado direito (o 1000). Se isso acontecer, ele acharia $x = 200$ (justificando também a letra B por outro caminho). Corrija lembrando da gangorra: o que faz de um lado, faz do outro.
    \end{tcolorbox}
    \vspace{0.15cm}
    
    % --- CAIXA VERDE (Porto Seguro Didático do Professor) ---
    \begin{boxexplicacao}
    \textbf{Porto Seguro Didático (O "Pulo do Gato"):} Odeia trabalhar com frações em equações? O MMC é o seu "raio desintegrador". Basta descobrir o MMC de todo mundo e multiplicar a equação inteira por ele (sem dó). Como num passe de mágica, todos os denominadores serão cortados e você voltará a trabalhar apenas com números inteiros no "chão de fábrica".
    \end{boxexplicacao}
\fi

\par\vspace{\espacoQuestao}


\noindent\textcolor{CorLinha}{\rule{\linewidth}{0.5pt}} 
\par\vspace{\espacoPreQuestao}
\subsubsection*{5.} 
\begin{enunciadoLiteral}
Em computação, a resolução de equações é ensinada às máquinas por meio de fluxogramas lógicos, que são passos sequenciais organizados em blocos de decisão. 

Crie, em formato de texto, um fluxograma sequencial listando pelo menos três passos fundamentais obrigatórios para programar um computador a resolver corretamente a equação $\mfrac{3x}{4} - 2 = x - 5$.
\end{enunciadoLiteral}

\ifgabarito
    % --- CAIXA LARANJA (Resolução e Tutor) ---
    \begin{tcolorbox}[colback=CorTitUm!5, colframe=CorTitUm, boxrule=1pt, arc=4pt, left=6pt, right=6pt, top=6pt, bottom=6pt]
    \small\textbf{\textcolor{CorTitUm}{Roteiro do Professor e Tutor IA:}}\par\vspace{0.1cm}
    \color{CinzaEscuro}
    \textbf{1. Ancoragem (Abordagem Socrática):} "Se você fosse um robô que só executa ordens simples, qual seria a primeira sujeira que você limparia desta equação para deixá-la mais fácil de ler?"\par\vspace{0.1cm}
    
    \textbf{2. Micro-passos da Resolução (Expectativa do Fluxograma):}\\
    O aluno deve descrever verbalmente ou em esquema a seguinte sequência (não os cálculos, mas os COMANDOS):\\
    \textit{Passo 1 (Remoção de Frações):} Multiplicar toda a equação pelo denominador comum (no caso, 4) para eliminar as frações.\\
    \textit{Passo 2 (Organização):} Aplicar o princípio aditivo, movendo todos os termos que contêm a variável \textbf{x} para o primeiro membro e as constantes numéricas para o segundo membro, invertendo as operações.\\
    \textit{Passo 3 (Isolamento Final):} Aplicar o princípio multiplicativo, dividindo ambos os lados da equação pelo coeficiente numérico que acompanha o \textbf{x}, garantindo o conjunto solução.\par\vspace{0.1cm}
    
    \textbf{3. Mapeamento de Armadilhas (Alerta Tutor):} O aluno pode querer resolver a equação no caderno e entregar apenas $x = 12$ como resposta. Não aceite. O comando cognitivo da questão é traduzir O MÉTODO em linguagem algorítmica, provando a metacognição.
    
    % MAPA CONCEITUAL TIKZ
    \begin{center}
    \begin{adjustbox}{trim=0cm 0cm 0cm 0cm, clip, max width=0.85\linewidth}
    \begin{tikzpicture}[node distance=1.5cm,
        box/.style={rectangle, draw=CinzaEscuro, rounded corners=3pt, fill=white, drop shadow={opacity=0.15, shadow xshift=1pt, shadow yshift=-1pt}, text width=5cm, align=center, minimum height=1cm, font=\small},
        seta/.style={->, thick, CinzaEscuro}]
        
        \node[box, fill=CorTitUm!10, draw=CorTitUm, font=\bfseries] (start) {INÍCIO};
        \node[box, below=0.6cm of start] (passo1) {Multiplicar todos os termos por 4 (eliminar fração)};
        \node[box, below=0.6cm of passo1] (passo2) {Agrupar variáveis de um lado e números do outro};
        \node[box, below=0.6cm of passo2] (passo3) {Dividir pelo coeficiente de x};
        \node[box, fill=CorNotaBorda!10, draw=CorNotaBorda, font=\bfseries, below=0.6cm of passo3] (end) {FIM: Conjunto Solução S};

        \draw[seta] (start) -- (passo1);
        \draw[seta] (passo1) -- (passo2);
        \draw[seta] (passo2) -- (passo3);
        \draw[seta] (passo3) -- (end);
    \end{tikzpicture}
    \end{adjustbox}
    \end{center}
    
    \end{tcolorbox}
    \vspace{0.15cm}
    
    % --- CAIXA VERDE (Porto Seguro Didático do Professor) ---
    \begin{boxexplicacao}
    \textbf{Porto Seguro Didático (O "Pulo do Gato"):} Ensinar matemática não é só fazer continha, é criar sistemas lógicos na cabeça. Quando o aluno consegue explicar os passos de como se faz algo (um tutorial), ele prova que dominou a "receita do bolo" e nunca mais esquecerá.
    \end{boxexplicacao}
\fi

\par\vspace{\espacoQuestao}


\noindent\textcolor{CorLinha}{\rule{\linewidth}{0.5pt}} 
\par\vspace{\espacoPreQuestao}
\subsubsection*{6.} 
\begin{enunciadoLiteral}
Considere o esquema a seguir representando a planta baixa de uma quadra poliesportiva de geometria retangular. O projeto informa que o contorno total (perímetro) pintado no chão corresponde a 108 metros lineares.

\begin{center}
% ==========================================
% CONTROLE INDIVIDUAL DESTA IMAGEM:
% 1. CORTAR BORDAS: trim=0.5cm 0cm 0.5cm 0cm
% 2. TAMANHO: \larguraTikz por um valor
% ==========================================
\begin{adjustbox}{trim=0cm 0cm 0cm 0cm, clip, max width=0.65\linewidth}
\begin{tikzpicture}
% --- APLICANDO DROP SHADOW E CORES ---
\definecolor{quadraFundo}{HTML}{E67E22}
\definecolor{quadraLinha}{HTML}{ECF0F1}

% LAYER 1: Grama/Piso externo (Opcional contextual)
\fill[gray!10, drop shadow={opacity=0.1, shadow xshift=1pt, shadow yshift=-2pt}, rounded corners=3pt] (-1,-1) rectangle (6,4);

% LAYER 2: A Quadra (Retângulo Principal)
\fill[quadraFundo, rounded corners=1pt] (0,0) rectangle (5,3);
\draw[quadraLinha, thick] (0,0) rectangle (5,3);

% LAYER 3: Detalhes da Quadra (Círculo central e linhas)
\draw[quadraLinha, thick] (2.5,0) -- (2.5,3);
\draw[quadraLinha, thick] (2.5,1.5) circle (0.6);
\draw[quadraLinha, thick] (0,1.0) rectangle (0.8,2.0);
\draw[quadraLinha, thick] (5,1.0) rectangle (4.2,2.0);

% LAYER 4: Cotas e Expressões Algébricas
\draw[|<->|, thick, Cinzablack] (0,-0.4) -- (5,-0.4) node[midway, fill=gray!10, font=\bfseries, text=black, inner sep=2pt] {3x - 4};
\draw[|<->|, thick, Cinzablack] (-0.4,0) -- (-0.4,3) node[midway, fill=gray!10, font=\bfseries, text=black, inner sep=2pt, rotate=90] {x + 10};

\end{tikzpicture}
\end{adjustbox}
\end{center}

Traduza as informações visuais da quadra para uma equação algébrica e determine o valor da constante \textbf{x}, que viabiliza a construção correta do projeto.
\end{enunciadoLiteral}

\ifgabarito
    % --- CAIXA LARANJA (Resolução e Tutor) ---
    \begin{tcolorbox}[colback=CorTitUm!5, colframe=CorTitUm, boxrule=1pt, arc=4pt, left=6pt, right=6pt, top=6pt, bottom=6pt]
    \small\textbf{\textcolor{CorTitUm}{Roteiro do Professor e Tutor IA:}}\par\vspace{0.1cm}
    \color{CinzaEscuro}
    \textbf{1. Ancoragem (Abordagem Socrática):} "Se olharmos para o desenho, temos dados em apenas dois lados. Mas o que é o perímetro de um retângulo? Quantos lados ele tem e como são os lados opostos?"\par\vspace{0.1cm}
    
    \textbf{2. Micro-passos da Resolução:}\\
    \textit{Passo 1:} O perímetro é a soma de todos os 4 lados. Como é um retângulo, os lados opostos são paralelos e iguais. Portanto, temos duas larguras e dois comprimentos.\\
    \textit{Passo 2:} Montagem estrutural $\rightarrow$ $2 \cdot (\text{comprimento}) + 2 \cdot (\text{largura}) = \text{Perímetro}$.\\
    \textit{Passo 3:} Substituição algébrica $\rightarrow$ $2 \cdot (3x - 4) + 2 \cdot (x + 10) = 108$.\\
    \textit{Passo 4:} Aplicação da distributiva $\rightarrow$ $6x - 8 + 2x + 20 = 108$.\\
    \textit{Passo 5:} Agrupamento $\rightarrow$ $8x + 12 = 108$.\\
    \textit{Passo 6:} Isolamento e solução $\rightarrow$ $8x = 108 - 12 \Rightarrow 8x = 96 \Rightarrow x = 12$.\par\vspace{0.1cm}
    
    \textbf{3. Mapeamento de Armadilhas (Alerta Tutor):} O erro imperdoável e mais frequente aqui é o aluno somar apenas os dois lados que aparecem escritos na imagem, gerando a equação $3x - 4 + x + 10 = 108$. Isso daria $x = 25,5$. Faça o aluno preencher a lápis as expressões nos quatro lados da figura antes de somar.
    \end{tcolorbox}
    \vspace{0.15cm}
    
    % --- CAIXA VERDE (Porto Seguro Didático do Professor) ---
    \begin{boxexplicacao}
    \textbf{Porto Seguro Didático (O "Pulo do Gato"):} Na geometria conectada à álgebra, o "invisível" também pesa. Nunca construa uma equação de perímetro sem antes dar a volta completa na figura apontando com o dedo. A palavra perímetro vem do grego (peri = em torno, metron = medida), ou seja, é a cerca completa da fazenda, não apenas duas cercas.
    \end{boxexplicacao}
\fi

\par\vspace{\espacoQuestao}


\noindent\textcolor{CorLinha}{\rule{\linewidth}{0.5pt}} 
\par\vspace{\espacoPreQuestao}
\subsubsection*{7.} 
\begin{enunciadoLiteral}
Na álgebra do primeiro grau, o princípio da equivalência dita as regras de segurança sobre o que podemos ou não fazer com os lados de uma igualdade sem que a balança perca o equilíbrio. Analise de forma minuciosa as afirmações matemáticas abaixo:

\begin{enumerate}[label=\Roman*., itemsep=\espacoAlt]
\item Ao adicionarmos ou subtrairmos o mesmo valor numérico simultaneamente nos dois membros de uma equação, suas raízes (conjunto solução) permanecem inalteradas.
\item Se elevarmos ao quadrado apenas o primeiro membro de uma equação do 1º grau, a balança se manterá equilibrada desde que o segundo membro seja multiplicado por zero.
\item A equação $2x = 10$ e a equação $6x = 30$ são consideradas equações equivalentes, pois possuem exatamente o mesmo conjunto solução.
\end{enumerate}

É matematicamente rigoroso e verdadeiro o que se afirma em:

\begin{enumerate}[label=\alph*), itemsep=\espacoAlt]
\item I e II, apenas.
\item I e III, apenas.
\item II e III, apenas.
\item I, II e III.
\end{enumerate}
\end{enunciadoLiteral}

\ifgabarito
    % --- CAIXA LARANJA (Resolução e Tutor) ---
    \begin{tcolorbox}[colback=CorTitUm!5, colframe=CorTitUm, boxrule=1pt, arc=4pt, left=6pt, right=6pt, top=6pt, bottom=6pt]
    \small\textbf{\textcolor{CorTitUm}{Roteiro do Professor e Tutor IA:}}\par\vspace{0.1cm}
    \color{CinzaEscuro}
    \textbf{1. Ancoragem (Abordagem Socrática):} "O que define se duas equações, mesmo sendo visualmente diferentes, são 'gêmeas de resultado'? Se fizermos algo de um lado e não do outro, a balança sobrevive?"\par\vspace{0.1cm}
    
    \textbf{2. Micro-passos da Resolução:}\\
    \textit{Passo 1:} Análise da I (VERDADEIRA). Este é o Princípio Aditivo. Somar ou subtrair o mesmo valor dos dois lados mantém o equilíbrio elástico perfeito.\\
    \textit{Passo 2:} Análise da II (FALSA). A afirmação sugere fazer operações completamente distintas nos membros (elevar ao quadrado de um lado, multiplicar por zero de outro). Isso destrói o Princípio da Equivalência, e multiplicar por zero "apaga" a informação do segundo membro, corrompendo a equação fatalmente.\\
    \textit{Passo 3:} Análise da III (VERDADEIRA). Equações equivalentes são aquelas que admitem a mesma raiz. Na primeira, $2x = 10 \Rightarrow x = 5$. Na segunda, $6x = 30 \Rightarrow x = 5$. Ambas partilham a solução $\{5\}$. O gabarito é \textbf{Letra B}.\par\vspace{0.1cm}
    
    \textbf{3. Mapeamento de Armadilhas (Alerta Tutor):} Alunos têm dificuldade com a palavra "equivalentes" na afirmativa III, achando que as equações precisam ser escritas com os mesmos algarismos. Peça para eles resolverem as duas rapidamente.
    \end{tcolorbox}
    \vspace{0.15cm}
    
    % --- CAIXA VERDE (Porto Seguro Didático do Professor) ---
    \begin{boxexplicacao}
    \textbf{Porto Seguro Didático (O "Pulo do Gato"):} Recapitulando: a palavra "equivalente" na álgebra não significa "igual no visual". Significa "mesmo destino, mesmos resultados". Um real é equivalente a quatro moedas de 25 centavos. Visualmente diferentes, mas o poder de compra (a raiz da equação) é exatamente o mesmo.
    \end{boxexplicacao}
\fi

\par\vspace{\espacoQuestao}


\noindent\textcolor{CorLinha}{\rule{\linewidth}{0.5pt}} 
\par\vspace{\espacoPreQuestao}
\subsubsection*{8.} 
\begin{enunciadoLiteral}
(ENEM - Adaptada) Uma empresa de logística analisa o custo para alugar pequenos furgões de carga com duas locadoras diferentes. A composição dos custos mensais funciona da seguinte forma:

\begin{itemize}[itemsep=\espacoAlt]
\item \textbf{Locadora Rápida:} Cobra uma taxa fixa de manutenção mensal de R\$ 1.200,00 mais R\$ 0,50 por quilômetro rodado.
\item \textbf{Locadora Segura:} Cobra uma taxa fixa de manutenção mensal de R\$ 900,00 mais R\$ 0,75 por quilômetro rodado.
\end{itemize}

O gerente de logística deseja saber o ponto de equilíbrio financeiro. Modelando essa situação por meio de uma equação do 1º grau, descubra a quilometragem exata (\textbf{k}) que a empresa deve rodar no mês para que o custo de aluguel seja rigorosamente o mesmo em qualquer uma das locadoras.
\end{enunciadoLiteral}

\ifgabarito
    % --- CAIXA LARANJA (Resolução e Tutor) ---
    \begin{tcolorbox}[colback=CorTitUm!5, colframe=CorTitUm, boxrule=1pt, arc=4pt, left=6pt, right=6pt, top=6pt, bottom=6pt]
    \small\textbf{\textcolor{CorTitUm}{Roteiro do Professor e Tutor IA:}}\par\vspace{0.1cm}
    \color{CinzaEscuro}
    \textbf{1. Ancoragem (Abordagem Socrática):} "No texto, a palavra-chave é 'mesmo custo'. Em linguagem matemática, qual símbolo usamos para representar coisas iguais? Vamos escrever uma 'receita' para o preço da Rápida e outra para a Segura e cruzá-las."\par\vspace{0.1cm}
    
    \textbf{2. Micro-passos da Resolução:}\\
    \textit{Passo 1:} Modelar o Custo da Rápida $\rightarrow$ $C_R = 1200 + 0,50k$.\\
    \textit{Passo 2:} Modelar o Custo da Segura $\rightarrow$ $C_S = 900 + 0,75k$.\\
    \textit{Passo 3:} Igualar as duas equações (o "ponto de equilíbrio") $\rightarrow$ $1200 + 0,50k = 900 + 0,75k$.\\
    \textit{Passo 4:} Agrupar os termos com \textbf{k} de um lado e valores fixos do outro $\rightarrow$ $1200 - 900 = 0,75k - 0,50k$.\\
    \textit{Passo 5:} Operacionalizar com decimais $\rightarrow$ $300 = 0,25k$.\\
    \textit{Passo 6:} Isolar o \textbf{k} $\rightarrow$ $k = \mfrac{300}{0,25}$.\\
    \textit{Passo 7:} Como dividir por 0,25 (um quarto) é o mesmo que multiplicar por 4, temos $k = 1200$ quilômetros.\par\vspace{0.1cm}
    
    \textbf{3. Mapeamento de Armadilhas (Alerta Tutor):} O aluno costuma travar na divisão final por número decimal ($\mfrac{300}{0,25}$). Se isso acontecer, sugira multiplicar tanto o numerador quanto o denominador por 100 para eliminar a vírgula ($\mfrac{30000}{25}$), destravando a operação braçal.
    \end{tcolorbox}
    \vspace{0.15cm}
    
    % --- CAIXA VERDE (Porto Seguro Didático do Professor) ---
    \begin{boxexplicacao}
    \textbf{Porto Seguro Didático (O "Pulo do Gato"):} Problemas de "ponto de equilíbrio" (comparar plano de celular, táxi, aluguel) são figurinhas carimbadas do ENEM. O macete estrutural é sempre o mesmo: descubra o "Fixo 1" e "Variável 1" e coloque de um lado do sinal de igual ($=$). Descubra o "Fixo 2" e "Variável 2" e coloque do outro lado. 
    \end{boxexplicacao}
\fi

\par\vspace{\espacoQuestao}


\noindent\textcolor{CorLinha}{\rule{\linewidth}{0.5pt}} 
\par\vspace{\espacoPreQuestao}
\subsubsection*{9.} 
\begin{enunciadoLiteral}
Uma caixa d'água municipal começou a apresentar vazamentos. Um fiscal da prefeitura fez a medição técnica e registrou no relatório de avarias que a água que restou no reservatório equivale a exatamente:

"A metade da capacidade total do tanque (\textbf{V}), somada à quarta parte da mesma capacidade total, resultando em exatos 180 mil litros de volume retido."

Monte a equação matemática fracionária que retrata a leitura do fiscal e resolva-a para descobrir qual é a capacidade total original \textbf{V} desse reservatório antes do vazamento.
\end{enunciadoLiteral}

\ifgabarito
    % --- CAIXA LARANJA (Resolução e Tutor) ---
    \begin{tcolorbox}[colback=CorTitUm!5, colframe=CorTitUm, boxrule=1pt, arc=4pt, left=6pt, right=6pt, top=6pt, bottom=6pt]
    \small\textbf{\textcolor{CorTitUm}{Roteiro do Professor e Tutor IA:}}\par\vspace{0.1cm}
    \color{CinzaEscuro}
    \textbf{1. Ancoragem (Abordagem Socrática):} "Como traduzimos 'a metade' e 'a quarta parte' em frações usando a letra V? E o que significa matematicamente a palavra 'resultando' no relatório?"\par\vspace{0.1cm}
    
    \textbf{2. Micro-passos da Resolução:}\\
    \textit{Passo 1:} Tradução da linguagem corrente para algébrica. "Metade" $\rightarrow$ $\mfrac{V}{2}$. "Quarta parte" $\rightarrow$ $\mfrac{V}{4}$.\\
    \textit{Passo 2:} Montagem da Equação $\rightarrow$ $\mfrac{V}{2} + \mfrac{V}{4} = 180$.\\
    \textit{Passo 3:} Tirar o MMC entre os denominadores 2, 4 e 1. O MMC é 4.\\
    \textit{Passo 4:} Multiplicar a equação inteira por 4 $\rightarrow$ $4 \cdot (\mfrac{V}{2}) + 4 \cdot (\mfrac{V}{4}) = 4 \cdot (180)$.\\
    \textit{Passo 5:} Redução ao nível do chão $\rightarrow$ $2V + 1V = 720$.\\
    \textit{Passo 6:} Agrupamento $\rightarrow$ $3V = 720$.\\
    \textit{Passo 7:} Solução $\rightarrow$ $V = \mfrac{720}{3} \Rightarrow V = 240$ mil litros.\par\vspace{0.1cm}
    
    \textbf{3. Mapeamento de Armadilhas (Alerta Tutor):} O aluno pode tentar fugir da equação montando um desenho ou lógica mental. Embora válido no dia a dia, force a escrita algébrica, pois a BNCC e as matrizes exigem a "tradução" do texto para a linguagem matemática simbólica.
    \end{tcolorbox}
    \vspace{0.15cm}
    
    % --- CAIXA VERDE (Porto Seguro Didático do Professor) ---
    \begin{boxexplicacao}
    \textbf{Porto Seguro Didático (O "Pulo do Gato"):} Recapitulando a regra do MMC: ao multiplicar a equação inteira pelo MMC, você está aplicando o Princípio Multiplicativo da balança de forma inteligente e matando o "bicho-papão" das frações no mesmo golpe.
    \end{boxexplicacao}
\fi

\par\vspace{\espacoQuestao}


\noindent\textcolor{CorLinha}{\rule{\linewidth}{0.5pt}}
\par\vspace{\espacoPreQuestao}
\subsubsection*{10.} 
\begin{enunciadoLiteral}
Numa fábrica, um braço robótico foi programado para colocar caixas em uma esteira pesadora. O manual da esteira apresenta um alerta em formato de equação para evitar sobrecarga (onde \textbf{c} é a massa da caixa):

$ 3(c - 2) + 2c = 4c + 10 $

Um técnico iniciante afirmou que a esteira vai suportar com segurança caixas de até \textbf{20 kg}. Usando manipulação algébrica e resolução de equações como base argumentativa, faça uma refutação analítica, provando matematicamente por que o técnico está perigosamente errado e qual é, de fato, a massa da caixa descrita na equação da máquina.
\end{enunciadoLiteral}

\ifgabarito
    % --- CAIXA LARANJA (Resolução e Tutor) ---
    \begin{tcolorbox}[colback=CorTitUm!5, colframe=CorTitUm, boxrule=1pt, arc=4pt, left=6pt, right=6pt, top=6pt, bottom=6pt]
    \small\textbf{\textcolor{CorTitUm}{Roteiro do Professor e Tutor IA:}}\par\vspace{0.1cm}
    \color{CinzaEscuro}
    \textbf{1. Ancoragem (Abordagem Socrática):} "O técnico deduziu algo. Mas como provamos que ele está errado sem basear nossa opinião em 'achismos'? Precisamos resolver a equação para encontrar a verdade mecânica."\par\vspace{0.1cm}
    
    \textbf{2. Micro-passos da Resolução:}\\
    \textit{Passo 1:} Iniciar a resolução algébrica distribuindo o 3 no primeiro membro $\rightarrow$ $3c - 6 + 2c = 4c + 10$.\\
    \textit{Passo 2:} Agrupar os termos semelhantes do lado esquerdo antes de mudar as peças de lugar $\rightarrow$ $5c - 6 = 4c + 10$.\\
    \textit{Passo 3:} Transportar e inverter sinais (Princípio Aditivo). Trazer o $4c$ para a esquerda e o $-6$ para a direita $\rightarrow$ $5c - 4c = 10 + 6$.\\
    \textit{Passo 4:} Redução final. Chega-se diretamente ao conjunto solução $\rightarrow$ $1c = 16 \Rightarrow c = 16 \text{ kg}$.\\
    \textit{Passo 5:} A refutação argumentativa. O técnico está errado porque o modelo matemático do limite da máquina exige que o equilíbrio de tensões ocorra exatamente em caixas de $16 \text{ kg}$. Colocar $20 \text{ kg}$ causaria sobrecarga e rompimento estrutural.\par\vspace{0.1cm}
    
    \textbf{3. Mapeamento de Armadilhas (Alerta Tutor):} O aluno que erra a distributiva do sinal ($3 \cdot -2$) gera um $+6$, montando $5c + 6 = 4c + 10$, o que daria $c = 4$. Reforce a revisão rigorosa do "chuveirinho" com sinais negativos, pois um erro desses num contexto prático quebra a máquina.
    \end{tcolorbox}
    \vspace{0.15cm}
    
    % --- CAIXA VERDE (Porto Seguro Didático do Professor) ---
    \begin{boxexplicacao}
    \textbf{Porto Seguro Didático (O "Pulo do Gato"):} Na matemática do mundo real (na física, na engenharia e na programação), a equação não é apenas um "joguinho de descobrir o X para tirar dez na prova". A equação é a alma do projeto. Errar o "X" significa que a ponte cai, a esteira quebra e o foguete desvia. A refutação algébrica é a prova pericial do engenheiro!
    \end{boxexplicacao}
\fi

\par\vspace{\espacoQuestao}
\subsection*{Capítulo 3: Problemas com Equações}
\vspace{-0.2cm}\noindent\rule{\linewidth}{1.5pt} 
\par\vspace{\espacoPosTitulo}
\subsubsection*{11.} 
\begin{enunciadoLiteral}
O professor de Matemática propôs um desafio lógico para a turma do 7º Ano. Ele disse: "Se vocês pegarem a minha idade atual, multiplicarem por três e, em seguida, subtraírem 15 anos desse resultado, encontrarão exatamente a idade do meu avô, que tem 84 anos". 

Traduza o desafio do professor para a linguagem algébrica e determine a idade atual dele.
\end{enunciadoLiteral}

\ifgabarito
    % --- CAIXA LARANJA (Resolução e Tutor) ---
    \begin{tcolorbox}[colback=CorTitUm!5, colframe=CorTitUm, boxrule=1pt, arc=4pt, left=6pt, right=6pt, top=6pt, bottom=6pt]
    \small\textbf{\textcolor{CorTitUm}{Roteiro do Professor e Tutor IA:}}\par\vspace{0.1cm}
    \color{CinzaEscuro}
    \textbf{1. Ancoragem (Abordagem Socrática):} "O que não conhecemos nessa história? A idade do professor. Vamos chamá-la de uma letra, por exemplo, \textbf{x}. Como escrevemos 'multiplicar por três' e 'subtrair 15' usando o \textbf{x}?"\par\vspace{0.1cm}
    
    \textbf{2. Micro-passos da Resolução:}\\
    \textit{Passo 1:} Definir a variável: Idade do professor $= x$.\\
    \textit{Passo 2:} Tradução algébrica: "multiplicar por três" ($3x$), "subtrair 15" ($- 15$), "exatamente a idade do avô" ($= 84$). Logo: $3x - 15 = 84$.\\
    \textit{Passo 3:} Princípio aditivo (isolar o \textbf{x}): $3x = 84 + 15 \Rightarrow 3x = 99$.\\
    \textit{Passo 4:} Princípio multiplicativo: $x = \mfrac{99}{3} \Rightarrow x = 33$ anos.\par\vspace{0.1cm}
    
    \textbf{3. Mapeamento de Armadilhas (Alerta Tutor):} O aluno pode errar a montagem escrevendo $15 - 3x = 84$, invertendo a ordem da subtração. Lembre-o de que a frase "subtraírem 15 anos desse resultado" indica que o 15 é quem está sendo retirado, portanto o sinal de menos acompanha o 15.
    \end{tcolorbox}
    \vspace{0.15cm}
    
    % --- CAIXA VERDE (Porto Seguro Didático do Professor) ---
    \begin{boxexplicacao}
    \textbf{Porto Seguro Didático (O "Pulo do Gato"):} Resolver problemas em matemática é como ser um tradutor de idiomas. Você precisa ler a frase em "Português" e escrever embaixo a exata mesma frase no idioma "Álgebra". Palavras como 'dobro', 'triplo' e 'metade' viram multiplicações e frações. Palavras como 'resulta', 'dá' e 'equivale' viram o sinal de igualdade ($=$).
    \end{boxexplicacao}
\fi

\par\vspace{\espacoQuestao}


\noindent\textcolor{CorLinha}{\rule{\linewidth}{0.5pt}} 
\par\vspace{\espacoPreQuestao}
\subsubsection*{12.} 
\begin{enunciadoLiteral}
Um marceneiro está cortando uma longa ripa de madeira para construir a moldura de um quadro. Ele fez dois cortes, dividindo a peça inteira em três segmentos, cujos comprimentos estão descritos algebricamente no esquema abaixo.

\begin{center}
% ==========================================
% CONTROLE INDIVIDUAL DESTA IMAGEM:
% 1. CORTAR BORDAS: trim=0cm 0cm 0cm 0cm
% 2. TAMANHO: \larguraTikz por um valor
% ==========================================
\begin{adjustbox}{trim=0cm 0cm 0cm 0cm, clip, max width=0.85\linewidth}
\begin{tikzpicture}
% --- APLICANDO DROP SHADOW E CORES ---
\definecolor{madeiraBase}{HTML}{D35400}
\definecolor{madeiraClara}{HTML}{E67E22}
\definecolor{corteCor}{HTML}{34495E}

% LAYER 1: Pedaço 1
\fill[madeiraClara, drop shadow={opacity=0.3, shadow xshift=1pt, shadow yshift=-2pt}, rounded corners=1pt] (0,0) rectangle (2.5,0.6);
\draw[corteCor, thick, rounded corners=1pt] (0,0) rectangle (2.5,0.6);
\node[font=\bfseries, text=black] at (1.25, 0.3) {x};

% LAYER 2: Pedaço 2
\fill[madeiraClara, drop shadow={opacity=0.3, shadow xshift=1pt, shadow yshift=-2pt}, rounded corners=1pt] (2.7,0) rectangle (6.2,0.6);
\draw[corteCor, thick, rounded corners=1pt] (2.7,0) rectangle (6.2,0.6);
\node[font=\bfseries, text=black] at (4.45, 0.3) {x + 10};

% LAYER 3: Pedaço 3
\fill[madeiraClara, drop shadow={opacity=0.3, shadow xshift=1pt, shadow yshift=-2pt}, rounded corners=1pt] (6.4,0) rectangle (10,0.6);
\draw[corteCor, thick, rounded corners=1pt] (6.4,0) rectangle (10,0.6);
\node[font=\bfseries, text=black] at (8.2, 0.3) {2x - 6};

% LAYER 4: Linhas de Cota Total
\draw[|<->|, thick, Cinzablack] (0,-0.5) -- (10,-0.5) node[midway, fill=white, font=\bfseries, inner sep=2pt] {Comprimento Total = 124 cm};
\end{tikzpicture}
\end{adjustbox}
\end{center}

Com base na interpretação visual e algébrica, calcule o valor da variável \textbf{x} e responda: qual é, em centímetros, a medida do maior pedaço cortado pelo marceneiro?
\end{enunciadoLiteral}

\ifgabarito
    % --- CAIXA LARANJA (Resolução e Tutor) ---
    \begin{tcolorbox}[colback=CorTitUm!5, colframe=CorTitUm, boxrule=1pt, arc=4pt, left=6pt, right=6pt, top=6pt, bottom=6pt]
    \small\textbf{\textcolor{CorTitUm}{Roteiro do Professor e Tutor IA:}}\par\vspace{0.1cm}
    \color{CinzaEscuro}
    \textbf{1. Ancoragem (Abordagem Socrática):} "Se juntarmos (somarmos) os três pedaços de madeira, qual deve ser o comprimento total da ripa? O que a linha com as setas no desenho nos diz?"\par\vspace{0.1cm}
    
    \textbf{2. Micro-passos da Resolução:}\\
    \textit{Passo 1:} Montagem. A soma dos segmentos é igual ao total $\rightarrow$ $x + (x + 10) + (2x - 6) = 124$.\\
    \textit{Passo 2:} Agrupamento dos termos semelhantes $\rightarrow$ $x + x + 2x = 4x$ e $10 - 6 = +4$. Equação simplificada: $4x + 4 = 124$.\\
    \textit{Passo 3:} Resolução algébrica $\rightarrow$ $4x = 124 - 4 \Rightarrow 4x = 120 \Rightarrow x = 30$.\\
    \textit{Passo 4:} Substituição. O problema pede o maior pedaço, não apenas o \textbf{x}. Substituindo $x=30$ nos três pedaços:\\
    - Pedaço 1: $30 \text{ cm}$\\
    - Pedaço 2: $30 + 10 = 40 \text{ cm}$\\
    - Pedaço 3: $2(30) - 6 = 60 - 6 = 54 \text{ cm}$.\\
    Portanto, o maior pedaço mede $54 \text{ cm}$.\par\vspace{0.1cm}
    
    \textbf{3. Mapeamento de Armadilhas (Alerta Tutor):} A maior armadilha não é o cálculo, é a interpretação final. O aluno acha o valor $x = 30$, sublinha, comemora e passa para a próxima. O Tutor deve perguntar: "O texto pediu o valor de x ou o tamanho do maior pedaço da madeira?".
    \end{tcolorbox}
    \vspace{0.15cm}
    
    % --- CAIXA VERDE (Porto Seguro Didático do Professor) ---
    \begin{boxexplicacao}
    \textbf{Porto Seguro Didático (O "Pulo do Gato"):} Matemática não é corrida de 100 metros onde achar o "x" é a linha de chegada. Muitas vezes, achar o "x" é só comprar a passagem; você ainda precisa viajar. Sempre releia a última frase da pergunta antes de entregar a resposta final.
    \end{boxexplicacao}
\fi

\par\vspace{\espacoQuestao}


\noindent\textcolor{CorLinha}{\rule{\linewidth}{0.5pt}}
\par\vspace{\espacoPreQuestao}
\subsubsection*{13.} 
\begin{enunciadoLiteral}
Uma questão clássica de olimpíada de matemática dizia o seguinte: "A soma de dois números inteiros e consecutivos é igual a 87. Quais são esses números?".

Um aluno tentou resolver o problema e modelou a seguinte equação: $x + y = 87$. Em seguida, ele travou, pois percebeu que uma equação com duas letras diferentes não poderia ser resolvida diretamente com as ferramentas que possuía.

\begin{boxresponda}
\textbf{Responda:} Qual erro de tradução algébrica o aluno cometeu ao usar $x$ e $y$? Mostre a equação correta utilizando apenas uma variável e resolva o problema, encontrando os dois números.
\end{boxresponda}
\end{enunciadoLiteral}

\ifgabarito
    % --- CAIXA LARANJA (Resolução e Tutor) ---
    \begin{tcolorbox}[colback=CorTitUm!5, colframe=CorTitUm, boxrule=1pt, arc=4pt, left=6pt, right=6pt, top=6pt, bottom=6pt]
    \small\textbf{\textcolor{CorTitUm}{Roteiro do Professor e Tutor IA:}}\par\vspace{0.1cm}
    \color{CinzaEscuro}
    \textbf{1. Ancoragem (Abordagem Socrática):} "O que significa a palavra 'consecutivo'? Se o primeiro número fosse 10, qual seria o consecutivo dele? Como você fez para chegar do 10 para o 11 matematicamente?"\par\vspace{0.1cm}
    
    \textbf{2. Micro-passos da Resolução:}\\
    \textit{Passo 1:} Diagnóstico do erro. Ao usar $x$ e $y$, o aluno modelou "a soma de dois números quaisquer", ignorando a palavra-chave "consecutivos".\\
    \textit{Passo 2:} Correção. Números consecutivos possuem uma diferença exata de 1 unidade. Se o primeiro número é \textbf{x}, o próximo será obrigatoriamente \textbf{x + 1}.\\
    \textit{Passo 3:} Equação correta $\rightarrow$ $x + (x + 1) = 87$.\\
    \textit{Passo 4:} Resolução $\rightarrow$ $2x + 1 = 87 \Rightarrow 2x = 86 \Rightarrow x = 43$.\\
    \textit{Passo 5:} Resposta final. O primeiro número (\textbf{x}) é 43. O número consecutivo (\textbf{x + 1}) é 44.\par\vspace{0.1cm}
    
    \textbf{3. Mapeamento de Armadilhas (Alerta Tutor):} Os alunos tendem a resolver de cabeça, tentando chutar números próximos a 40 (ex: $40+47$). Não aceite tentativas de "tentativa e erro". Exija que o aluno escreva o modelo algébrico, pois essa é a habilidade avaliada no capítulo.
    \end{tcolorbox}
    \vspace{0.15cm}
    
    % --- CAIXA VERDE (Nota Pedagógica) ---
    \begin{boxexplicacao}
    \textbf{Justificativa Curricular (Raio-X da Questão):} Avalia a habilidade de traduzir relações de dependência (consecutivos, pares, ímpares) para a álgebra usando a mesma base variável, rompendo o vício de criar uma letra nova para cada elemento do texto.
    \end{boxexplicacao}
\fi

\par\vspace{\espacoQuestao}


\noindent\textcolor{CorLinha}{\rule{\linewidth}{0.5pt}} 
\par\vspace{\espacoPreQuestao}
\subsubsection*{14.} 
\begin{enunciadoLiteral}
Em um aplicativo de transporte, o custo de uma viagem é composto por um valor fixo de embarque (a "bandeirada") de R\$ 4,50, mais um valor variável de R\$ 1,20 para cada quilômetro rodado. 

Um passageiro pagou exatamente R\$ 28,50 por uma corrida. Sendo \textbf{k} a quantidade de quilômetros rodados, marque a alternativa que indica corretamente o tamanho do trajeto percorrido e, em seguida, apresente os cálculos provando por que a alternativa D é matematicamente falsa.

\begin{enumerate}[label=\alph*), itemsep=\espacoAlt]
\item $ k = 15 \text{ km} $
\item $ k = 18 \text{ km} $
\item $ k = 20 \text{ km} $
\item $ k = 27 \text{ km} $
\end{enumerate}
\end{enunciadoLiteral}

\ifgabarito
    % --- CAIXA LARANJA (Resolução e Tutor) ---
    \begin{tcolorbox}[colback=CorTitUm!5, colframe=CorTitUm, boxrule=1pt, arc=4pt, left=6pt, right=6pt, top=6pt, bottom=6pt]
    \small\textbf{\textcolor{CorTitUm}{Roteiro do Professor e Tutor IA:}}\par\vspace{0.1cm}
    \color{CinzaEscuro}
    \textbf{1. Ancoragem (Abordagem Socrática):} "O valor final da corrida depende de quê? A bandeirada muda se a corrida for longa ou curta? Como juntamos o valor fixo com o valor variável que depende de 'k'?"\par\vspace{0.1cm}
    
    \textbf{2. Micro-passos da Resolução:}\\
    \textit{Passo 1:} Modelagem. Custo Total = Fixo + (Variável $\cdot$ quilômetros) $\rightarrow$ $4,50 + 1,20k = 28,50$.\\
    \textit{Passo 2:} Isolamento. Mover o custo fixo $\rightarrow$ $1,20k = 28,50 - 4,50 \Rightarrow 1,20k = 24,00$.\\
    \textit{Passo 3:} Divisão. $k = \mfrac{24}{1,20}$. Para facilitar, multiplica-se numerador e denominador por 10: $k = \mfrac{240}{12} \Rightarrow k = 20 \text{ km}$. A alternativa correta é a \textbf{letra C}.\\
    \textit{Passo 4:} Refutação da Letra D. O erro da letra D ($k = 27$) ocorre quando o passageiro soma de forma equivocada os valores do enunciado ($28,50 - 1,50 \text{ (erro de conta)} = 27$), ou ignora totalmente o modelo e tenta chutar valores aleatórios que beiram os 30 km.\par\vspace{0.1cm}
    
    \textbf{3. Mapeamento de Armadilhas (Alerta Tutor):} Operar com decimais ($1,20$) costuma paralisar o aluno. Lembre-o de que o zero à direita após a vírgula não tem valor ($1,2$) e use o truque de transformar dinheiro em centavos se preferir ($450 + 120k = 2850$).
    \end{tcolorbox}
    \vspace{0.15cm}
    
    % --- CAIXA VERDE (Porto Seguro Didático do Professor) ---
    \begin{boxexplicacao}
    \textbf{Porto Seguro Didático (O "Pulo do Gato"):} Recapitulando a modelagem de negócios: Valor Fixo + (Valor Variável $\times$ Quantidade) = Total. Essa é a estrutura mestre de praticamente todas as contas de táxi, telefone, aluguel de carros e salário de vendedores!
    \end{boxexplicacao}
\fi

\par\vspace{\espacoQuestao}


\noindent\textcolor{CorLinha}{\rule{\linewidth}{0.5pt}} 
\par\vspace{\espacoPreQuestao}
\subsubsection*{15.} 
\begin{enunciadoLiteral}
A linguagem algébrica permite sintetizar frases inteiras em pequenas fórmulas precisas. Escreva, em formato de itens, a tradução matemática para cada uma das sentenças a seguir. Utilize sempre a letra \textbf{n} para representar o "número desconhecido".

1. "O dobro de um número, adicionado à sua terça parte."
2. "A metade do sucessor de um número."
3. "A diferença entre o quádruplo de um número e 15."
\end{enunciadoLiteral}

\ifgabarito
    % --- CAIXA LARANJA (Resolução e Tutor) ---
    \begin{tcolorbox}[colback=CorTitUm!5, colframe=CorTitUm, boxrule=1pt, arc=4pt, left=6pt, right=6pt, top=6pt, bottom=6pt]
    \small\textbf{\textcolor{CorTitUm}{Roteiro do Professor e Tutor IA:}}\par\vspace{0.1cm}
    \color{CinzaEscuro}
    \textbf{1. Ancoragem (Abordagem Socrática):} "Vamos por partes. Como se escreve 'dobro'? E 'terça parte'? Na frase 2, quem deve ser dividido pela metade: o número sozinho ou o sucessor dele inteiro?"\par\vspace{0.1cm}
    
    \textbf{2. Micro-passos da Resolução:}\\
    \textit{Tradução 1:} "Dobro" é $2n$. "Adicionado" é $+$. "Terça parte" é $\mfrac{n}{3}$. Expressão: $2n + \mfrac{n}{3}$.\\
    \textit{Tradução 2:} "Sucessor" de um número é $n + 1$. A "metade" deve abraçar todo o sucessor, exigindo uma fração global. Expressão: $\mfrac{n + 1}{2}$.\\
    \textit{Tradução 3:} "Diferença" é o sinal de $-$. "Quádruplo" é $4n$. Expressão: $4n - 15$.\par\vspace{0.1cm}
    
    \textbf{3. Mapeamento de Armadilhas (Alerta Tutor):} O perigo maior está na frase 2. O aluno costuma traduzir literalmente da esquerda para a direita e escrever $\mfrac{n}{2} + 1$. Isso é "a metade do número mais um", e não "a metade do sucessor". Force a leitura semântica do bloco antes de escrever.
    \end{tcolorbox}
    \vspace{0.15cm}
    
    % --- CAIXA VERDE (Porto Seguro Didático do Professor) ---
    \begin{boxexplicacao}
    \textbf{Porto Seguro Didático (O "Pulo do Gato"):} A língua portuguesa engana a matemática! A vírgula ou a ordem das palavras muda toda a conta. "O dobro de (um número mais cinco)" é $2 \cdot (x + 5)$. Já "O dobro de um número, mais cinco" é $2x + 5$. Muita atenção à sintaxe textual!
    \end{boxexplicacao}
\fi

\par\vspace{\espacoQuestao}


\noindent\textcolor{CorLinha}{\rule{\linewidth}{0.5pt}} 
\par\vspace{\espacoPreQuestao}
\subsubsection*{16.} 
\begin{enunciadoLiteral}
No layout arquitetônico de um laboratório, uma grande esteira de triagem eletrônica precisa ser instalada exatamente entre duas vigas mestras de sustentação. O espaço total disponível de parede a parede está cotado no projeto abaixo. A esteira é composta por três módulos idênticos de comprimento \textbf{x} e dois espaços de manutenção obrigatórios de \textbf{2,5 metros} cada.

\begin{center}
% ==========================================
% CONTROLE INDIVIDUAL DESTA IMAGEM:
% 1. CORTAR BORDAS: trim=0cm 0cm 0cm 0cm
% 2. TAMANHO: \larguraTikz por um valor
% ==========================================
\begin{adjustbox}{trim=0cm 0cm 0cm 0cm, clip, max width=0.9\linewidth}
\begin{tikzpicture}
% --- APLICANDO DROP SHADOW E CORES ---
\definecolor{vigaCor}{HTML}{7F8C8D}
\definecolor{esteiraCor}{HTML}{2980B9}
\definecolor{espacoCor}{HTML}{BDC3C7}

% LAYER 1: Vigas laterais
\fill[vigaCor, drop shadow={opacity=0.3, shadow xshift=1pt, shadow yshift=-2pt}] (0,0) rectangle (0.8,3);
\fill[vigaCor, drop shadow={opacity=0.3, shadow xshift=1pt, shadow yshift=-2pt}] (11.2,0) rectangle (12,3);
\node[font=\bfseries, text=white, rotate=90] at (0.4, 1.5) {Viga Esquerda};
\node[font=\bfseries, text=white, rotate=90] at (11.6, 1.5) {Viga Direita};

% LAYER 2: Módulos da Esteira
\fill[esteiraCor, rounded corners=2pt, drop shadow={opacity=0.2, shadow xshift=1pt, shadow yshift=-2pt}] (0.8,1) rectangle (3.3,1.8);
\node[font=\bfseries, text=white] at (2.05, 1.4) {Módulo x};

\fill[esteiraCor, rounded corners=2pt, drop shadow={opacity=0.2, shadow xshift=1pt, shadow yshift=-2pt}] (4.7,1) rectangle (7.2,1.8);
\node[font=\bfseries, text=white] at (5.95, 1.4) {Módulo x};

\fill[esteiraCor, rounded corners=2pt, drop shadow={opacity=0.2, shadow xshift=1pt, shadow yshift=-2pt}] (8.6,1) rectangle (11.1,1.8);
\node[font=\bfseries, text=white] at (9.85, 1.4) {Módulo x};

% LAYER 3: Espaços de Manutenção (Conectores)
\fill[espacoCor, drop shadow={opacity=0.2, shadow xshift=1pt, shadow yshift=-1pt}] (3.3,1.3) rectangle (4.7,1.5);
\node[font=\small\bfseries, text=black] at (4.0, 1.0) {2,5 m};

\fill[espacoCor, drop shadow={opacity=0.2, shadow xshift=1pt, shadow yshift=-1pt}] (7.2,1.3) rectangle (8.6,1.5);
\node[font=\small\bfseries, text=black] at (7.9, 1.0) {2,5 m};

% LAYER 4: Cota Total (Distância entre vigas)
\draw[|<->|, thick, Cinzablack] (0.8,-0.5) -- (11.2,-0.5) node[midway, fill=white, font=\bfseries, inner sep=2pt] {Comprimento Vão Livre = 26 metros};

\end{tikzpicture}
\end{adjustbox}
\end{center}

Considerando que a esteira coube perfeitamente e encosta nas duas vigas sem sobras, modele a equação que descreve esse sistema e descubra o comprimento unitário (\textbf{x}) de cada módulo.
\end{enunciadoLiteral}

\ifgabarito
    % --- CAIXA LARANJA (Resolução e Tutor) ---
    \begin{tcolorbox}[colback=CorTitUm!5, colframe=CorTitUm, boxrule=1pt, arc=4pt, left=6pt, right=6pt, top=6pt, bottom=6pt]
    \small\textbf{\textcolor{CorTitUm}{Roteiro do Professor e Tutor IA:}}\par\vspace{0.1cm}
    \color{CinzaEscuro}
    \textbf{1. Ancoragem (Abordagem Socrática):} "Quantos módulos com letra 'x' nós temos enfileirados? E quantos espaços de 2,5 metros existem entre eles? Juntando tudo, que distância eles precisam preencher?"\par\vspace{0.1cm}
    
    \textbf{2. Micro-passos da Resolução:}\\
    \textit{Passo 1:} Leitura visual da planta (Teste Cego). Temos 3 módulos \textbf{x} e 2 espaços de manutenção de $2,5 \text{ m}$.\\
    \textit{Passo 2:} Modelagem matemática $\rightarrow$ $x + 2,5 + x + 2,5 + x = 26$.\\
    \textit{Passo 3:} Redução dos termos $\rightarrow$ $3x + 5 = 26$.\\
    \textit{Passo 4:} Isolamento da incógnita $\rightarrow$ $3x = 26 - 5 \Rightarrow 3x = 21$.\\
    \textit{Passo 5:} Solução Final $\rightarrow$ $x = \mfrac{21}{3} \Rightarrow x = 7 \text{ metros}$.\par\vspace{0.1cm}
    
    \textbf{3. Mapeamento de Armadilhas (Alerta Tutor):} O aluno descuidado costuma ler "3 módulos e espaços de 2,5" e modelar $3x + 2,5 = 26$, esquecendo que o esquema exige DOIS espaços de manutenção para conectar três peças. Aponte para o desenho e mande-o contar com o dedo.
    \end{tcolorbox}
    \vspace{0.15cm}
    
    % --- CAIXA VERDE (Porto Seguro Didático do Professor) ---
    \begin{boxexplicacao}
    \textbf{Porto Seguro Didático (O "Pulo do Gato"):} Recapitulando a regra da geometria aplicada: o todo é igual à soma exata das suas partes. A álgebra é só a ferramenta que usamos para costurar esses retalhos em uma linha contínua!
    \end{boxexplicacao}
\fi

\par\vspace{\espacoQuestao}


\noindent\textcolor{CorLinha}{\rule{\linewidth}{0.5pt}} 
\par\vspace{\espacoPreQuestao}
\subsubsection*{17.} 
\begin{enunciadoLiteral}
(Estilo VUNESP) Em um pátio de fiscalização de trânsito estão apreendidos apenas carros (veículos de 4 rodas) e motocicletas (veículos de 2 rodas). O inspetor de pátio anotou que há um total de 60 veículos estacionados. Durante a contagem, ele observou uma proporção rigorosa: o número de carros é exatamente o triplo do número de motocicletas.

Não é necessário saber o total de rodas para resolver o caso. Utilizando apenas uma equação do 1º grau (com uma única variável), determine a quantidade exata de motocicletas e de carros presentes no pátio.
\end{enunciadoLiteral}

\ifgabarito
    % --- CAIXA LARANJA (Resolução e Tutor) ---
    \begin{tcolorbox}[colback=CorTitUm!5, colframe=CorTitUm, boxrule=1pt, arc=4pt, left=6pt, right=6pt, top=6pt, bottom=6pt]
    \small\textbf{\textcolor{CorTitUm}{Roteiro do Professor e Tutor IA:}}\par\vspace{0.1cm}
    \color{CinzaEscuro}
    \textbf{1. Ancoragem (Abordagem Socrática):} "Se existem duas coisas desconhecidas (carros e motos), como usamos a frase 'o número de carros é o triplo das motos' para usar a mesma letra para os dois?"\par\vspace{0.1cm}
    
    \textbf{2. Micro-passos da Resolução:}\\
    \textit{Passo 1:} Definição da base algébrica. Como os carros dependem do número de motos, chamamos as motocicletas de \textbf{m}. Logo, os carros serão \textbf{3m}.\\
    \textit{Passo 2:} Montagem da Equação. O total de veículos soma 60. Equação: $m + 3m = 60$.\\
    \textit{Passo 3:} Resolução $\rightarrow$ $4m = 60 \Rightarrow m = \mfrac{60}{4} \Rightarrow m = 15$ motocicletas.\\
    \textit{Passo 4:} Determinar os carros. Como carros $= 3m$, calculamos $3 \cdot 15 = 45$ carros.\\
    \textit{Passo 5:} Prova real. $15 \text{ (motos)} + 45 \text{ (carros)} = 60 \text{ veículos}$. O cálculo bate.\par\vspace{0.1cm}
    
    \textbf{3. Mapeamento de Armadilhas (Alerta Tutor):} Alunos do 7º ano muitas vezes tentam forçar um Sistema de Equações ($C + M = 60$) e travam por não terem bagagem ainda. Direcione-os a escrever tudo em função de uma única variável. A pergunta "quem depende de quem" salva o problema.
    \end{tcolorbox}
    \vspace{0.15cm}
    
    % --- CAIXA VERDE (Porto Seguro Didático do Professor) ---
    \begin{boxexplicacao}
    \textbf{Porto Seguro Didático (O "Pulo do Gato"):} Quando um problema apresenta dois "desconhecidos" (carros e motos, João e Maria, adultos e crianças), mas te dá uma dica de proporção (triplo, metade, mais dez), você NÃO precisa de duas letras. Dê o "x" para o elo mais fraco (a base) e construa a expressão do outro dependendo desse "x". 
    \end{boxexplicacao}
\fi

\par\vspace{\espacoQuestao}


\noindent\textcolor{CorLinha}{\rule{\linewidth}{0.5pt}} 
\par\vspace{\espacoPreQuestao}
\subsubsection*{18.} 
\begin{enunciadoLiteral}
Uma rede de academias de ginástica oferece dois planos de mensalidade para seus clientes, dependendo de quantas aulas semanais o aluno fará no mês:

\begin{itemize}[itemsep=\espacoAlt]
\item \textbf{Plano Light:} Taxa de matrícula fixa anual de R\$ 120,00, mais R\$ 15,00 por aula frequentada.
\item \textbf{Plano Pro:} Taxa de matrícula fixa anual de R\$ 180,00, mais R\$ 10,00 por aula frequentada.
\end{itemize}

Seja \textbf{x} o número de aulas que um cliente frequenta ao longo do ano. Avalie as afirmações a seguir:

\begin{enumerate}[label=\Roman*., itemsep=\espacoAlt]
\item A equação que determina o ponto de equilíbrio (onde os dois planos custam o mesmo valor) é $120 + 15x = 180 + 10x$.
\item Resolvendo a equação do ponto de equilíbrio, descobre-se que ela ocorre exatamente quando o aluno frequenta 12 aulas.
\item Se um aluno pretende fazer 20 aulas no ano, o Plano Light sairá matematicamente mais barato que o Plano Pro.
\end{enumerate}

É correto o que se afirma em:

\begin{enumerate}[label=\alph*), itemsep=\espacoAlt]
\item I, apenas.
\item I e II, apenas.
\item II e III, apenas.
\item I, II e III.
\end{enumerate}
\end{enunciadoLiteral}

\ifgabarito
    % --- CAIXA LARANJA (Resolução e Tutor) ---
    \begin{tcolorbox}[colback=CorTitUm!5, colframe=CorTitUm, boxrule=1pt, arc=4pt, left=6pt, right=6pt, top=6pt, bottom=6pt]
    \small\textbf{\textcolor{CorTitUm}{Roteiro do Professor e Tutor IA:}}\par\vspace{0.1cm}
    \color{CinzaEscuro}
    \textbf{1. Ancoragem (Abordagem Socrática):} "Como testamos se um plano é mais vantajoso que o outro? Primeiro, achamos o momento exato em que há um 'empate' de preços (equilíbrio) e depois simulamos valores."\par\vspace{0.1cm}
    
    \textbf{2. Micro-passos da Resolução:}\\
    \textit{Passo 1:} Análise da I (VERDADEIRA). O modelo está perfeito. Plano Light $= 120 + 15x$. Plano Pro $= 180 + 10x$.\\
    \textit{Passo 2:} Análise da II (VERDADEIRA). Resolvendo a equação do equilíbrio: $15x - 10x = 180 - 120 \Rightarrow 5x = 60 \Rightarrow x = 12$ aulas.\\
    \textit{Passo 3:} Análise da III (FALSA). Testando para $x = 20$ aulas:\\
    Custo Light $= 120 + 15(20) = 120 + 300 = \text{R\$ } 420,00$.\\
    Custo Pro $= 180 + 10(20) = 180 + 200 = \text{R\$ } 380,00$.\\
    Logo, para 20 aulas, o Plano Pro é mais barato. A afirmativa III mentiu. Resposta: \textbf{Letra B}.\par\vspace{0.1cm}
    
    \textbf{3. Mapeamento de Armadilhas (Alerta Tutor):} Os alunos podem julgar a III como verdadeira no instinto, porque associam a palavra "Light" a algo mais barato, sem calcular matematicamente. Obrigue a testagem de hipóteses braçal.
    \end{tcolorbox}
    \vspace{0.15cm}
    
    % --- CAIXA VERDE (Porto Seguro Didático do Professor) ---
    \begin{boxexplicacao}
    \textbf{Porto Seguro Didático (O "Pulo do Gato"):} O ponto de equilíbrio é um divisor de águas! Se o empate deu 12 aulas, significa que abaixo de 12 aulas, vale a pena pagar a matrícula mais barata (Light). Mas se o aluno for assíduo e passar das 12 aulas, o custo variável mais alto do Light vai destruir a carteira dele, fazendo o Pro ser a escolha inteligente.
    \end{boxexplicacao}
\fi

\par\vspace{\espacoQuestao}


\noindent\textcolor{CorLinha}{\rule{\linewidth}{0.5pt}} 
\par\vspace{\espacoPreQuestao}
\subsubsection*{19.} 
\begin{enunciadoLiteral}
Em um problema contextualizado sobre descontos, constava o seguinte desafio: 
"Um par de tênis estava sendo vendido, em promoção, com R\$ 40,00 de desconto no caixa. Devido à promoção, o valor final a ser pago foi parcelado em 3 vezes iguais (sem juros) de R\$ 60,00 cada. Qual era o preço original desse tênis?"

Um aluno da turma modelou e resolveu a equação da seguinte forma:
$ \mfrac{x}{3} - 40 = 60 \rightarrow \mfrac{x}{3} = 100 \rightarrow x = 300 $ reais.

Ao chegar no caixa, se a etiqueta marcasse 300 reais, a matemática não bateria com as parcelas descritas.

\begin{boxresponda}
\textbf{Responda:} Qual foi o erro grotesco de modelagem do aluno? Mostre a montagem algébrica correta e encontre o real valor da etiqueta original (\textbf{x}).
\end{boxresponda}
\end{enunciadoLiteral}

\ifgabarito
    % --- CAIXA LARANJA (Resolução e Tutor) ---
    \begin{tcolorbox}[colback=CorTitUm!5, colframe=CorTitUm, boxrule=1pt, arc=4pt, left=6pt, right=6pt, top=6pt, bottom=6pt]
    \small\textbf{\textcolor{CorTitUm}{Roteiro do Professor e Tutor IA:}}\par\vspace{0.1cm}
    \color{CinzaEscuro}
    \textbf{1. Ancoragem (Abordagem Socrática):} "O que aconteceu primeiro no mundo real: a loja dividiu o valor em 3 vezes e depois deu o desconto em uma das parcelas, ou a loja deu o desconto no total e depois dividiu o que sobrou?"\par\vspace{0.1cm}
    
    \textbf{2. Micro-passos da Resolução:}\\
    \textit{Passo 1:} Diagnóstico do Erro. A equação do aluno $\mfrac{x}{3} - 40$ significa que ele pegou o preço total, dividiu em 3 parcelas e aplicou um desconto de 40 apenas na parcela. Isso subverte a cronologia comercial lógica.\\
    \textit{Passo 2:} Correção da Modelagem. O desconto de 40 acontece sobre o valor original inteiro (\textbf{x}). Todo esse bloco $(x - 40)$ é que sofre a divisão por 3. A equação correta é: $\mfrac{x - 40}{3} = 60$.\\
    \textit{Passo 3:} Resolução $\rightarrow$ O 3 que está dividindo tudo passa multiplicando: $x - 40 = 60 \cdot 3 \Rightarrow x - 40 = 180$.\\
    \textit{Passo 4:} Solução Final $\rightarrow$ $x = 180 + 40 \Rightarrow x = 220$ reais. O tênis custava R\$ 220,00.\par\vspace{0.1cm}
    
    \textbf{3. Mapeamento de Armadilhas (Alerta Tutor):} Frações globais requerem atenção. Tudo que está no numerador (acima da linha de fração) funciona como se estivesse entre parênteses inquebráveis. Você precisa lidar com o denominador (o 3) primeiro, antes de tocar no 40.
    \end{tcolorbox}
    \vspace{0.15cm}
    
    % --- CAIXA VERDE (Nota Pedagógica) ---
    \begin{boxexplicacao}
    \textbf{Justificativa Curricular (Raio-X da Questão):} Mede a capacidade do aluno de transitar da linguagem escrita para a modelagem algébrica respeitando a ordem cronológica e hierárquica das operações. Avalia o uso correto do traço de fração como agrupador sintático.
    \end{boxexplicacao}
\fi

\par\vspace{\espacoQuestao}


\noindent\textcolor{CorLinha}{\rule{\linewidth}{0.5pt}} 
\par\vspace{\espacoPreQuestao}
\subsubsection*{20.} 
\begin{enunciadoLiteral}
Uma prefeitura municipal aprovou a desapropriação de dois terrenos vizinhos para a construção de um pequeno posto de saúde. O projeto arquitetônico exige que os dois lotes tenham, matematicamente, áreas absolutamente iguais, embora tenham formatos diferentes.

\begin{center}
% ==========================================
% CONTROLE INDIVIDUAL DESTA IMAGEM:
% 1. CORTAR BORDAS: trim=0cm 0cm 0cm 0cm
% 2. TAMANHO: \larguraTikz por um valor
% ==========================================
\begin{adjustbox}{trim=0cm 0cm 0cm 0cm, clip, max width=0.85\linewidth}
\begin{tikzpicture}
% --- APLICANDO DROP SHADOW E CORES ---
\definecolor{terrenoUm}{HTML}{27AE60}
\definecolor{terrenoDois}{HTML}{2980B9}
\definecolor{asfalto}{HTML}{2F3640}

% LAYER 1: Rua / Asfalto
\fill[asfalto, drop shadow={opacity=0.3, shadow xshift=2pt, shadow yshift=-2pt}] (-1, -1) rectangle (12, 0);
\draw[white, dashed, thick] (-1, -0.5) -- (12, -0.5);
\node[font=\bfseries, text=white] at (5.5, -0.2) {Rua Principal};

% LAYER 2: Lote A
\fill[terrenoUm, drop shadow={opacity=0.2, shadow xshift=2pt, shadow yshift=-2pt}] (0,0) rectangle (4,5);
\draw[Cinzablack, thick] (0,0) rectangle (4,5);
\node[font=\large\bfseries, text=white] at (2, 2.5) {LOTE A};
\node[font=\bfseries, text=black] at (2, 5.3) {x};
\node[font=\bfseries, text=black, rotate=90] at (-0.3, 2.5) {5 metros};

% LAYER 3: Lote B
\fill[terrenoDois, drop shadow={opacity=0.2, shadow xshift=2pt, shadow yshift=-2pt}] (5,0) rectangle (11,4);
\draw[Cinzablack, thick] (5,0) rectangle (11,4);
\node[font=\large\bfseries, text=white] at (8, 2) {LOTE B};
\node[font=\bfseries, text=black] at (8, 4.3) {x + 2};
\node[font=\bfseries, text=black, rotate=90] at (11.3, 2) {4 metros};

\end{tikzpicture}
\end{adjustbox}
\end{center}

Lembrando que a área de um retângulo é calculada multiplicando-se a sua base pela sua altura, modele a equação mestre deste problema, descubra o valor da constante \textbf{x} e responda: qual é a área total de cada lote em metros quadrados?
\end{enunciadoLiteral}

\ifgabarito
    % --- CAIXA LARANJA (Resolução e Tutor) ---
    \begin{tcolorbox}[colback=CorTitUm!5, colframe=CorTitUm, boxrule=1pt, arc=4pt, left=6pt, right=6pt, top=6pt, bottom=6pt]
    \small\textbf{\textcolor{CorTitUm}{Roteiro do Professor e Tutor IA:}}\par\vspace{0.1cm}
    \color{CinzaEscuro}
    \textbf{1. Ancoragem (Abordagem Socrática):} "O texto nos dá uma ordem direta e absoluta: as áreas devem ser iguais. Como montamos a conta de área do Lote A? E a do Lote B, lembrando que a base dele tem dois termos?"\par\vspace{0.1cm}
    
    \textbf{2. Micro-passos da Resolução:}\\
    \textit{Passo 1:} Expressão da Área A $\rightarrow$ $\text{Base} \cdot \text{Altura} = x \cdot 5 = 5x$.\\
    \textit{Passo 2:} Expressão da Área B $\rightarrow$ $\text{Base} \cdot \text{Altura} = (x + 2) \cdot 4$. O parênteses é obrigatório porque o 4 multiplica tudo.\\
    \textit{Passo 3:} Equacionamento. Área A $=$ Área B $\rightarrow$ $5x = 4(x + 2)$.\\
    \textit{Passo 4:} Distributiva no Lote B $\rightarrow$ $5x = 4x + 8$.\\
    \textit{Passo 5:} Resolução $\rightarrow$ $5x - 4x = 8 \Rightarrow x = 8$.\\
    \textit{Passo 6:} Resposta Final. A pergunta não é apenas o valor de x, mas a Área do Lote. Substituindo $x=8$ no Lote A: Área $= 5 \cdot 8 = 40 \text{ m}^2$. (No Lote B seria $4 \cdot (8+2) = 4 \cdot 10 = 40 \text{ m}^2$, provando a igualdade).\par\vspace{0.1cm}
    
    \textbf{3. Mapeamento de Armadilhas (Alerta Tutor):} Omitir os parênteses no Lote B é o erro fatal mais comum (modelar $5x = x + 2 \cdot 4$, gerando $5x = x + 8 \Rightarrow 4x = 8 \Rightarrow x=2$). Vigie de perto a aplicação da propriedade distributiva quando houver bases compostas.
    \end{tcolorbox}
    \vspace{0.15cm}
    
    % --- CAIXA VERDE (Porto Seguro Didático do Professor) ---
    \begin{boxexplicacao}
    \textbf{Porto Seguro Didático (O "Pulo do Gato"):} Na geometria, sempre que o lado de uma figura for uma expressão algébrica composta (com sinal de mais ou menos, tipo $x+2$), essa expressão VIVE permanentemente dentro de parênteses invisíveis. Se ela for se chocar com qualquer outra operação (multiplicação de área, etc), os parênteses precisam aparecer para proteger o bloco todo!
    \end{boxexplicacao}
\fi

\par\vspace{\espacoQuestao}
\subsection*{Capítulo 4: Inequações do 1º Grau}
\vspace{-0.2cm}\noindent\rule{\linewidth}{1.5pt} 
\par\vspace{\espacoPosTitulo}
\subsubsection*{21.} 
\begin{enunciadoLiteral}
Uma desigualdade matemática difere de uma equação porque ela não busca um valor único e isolado, mas sim uma região inteira de números reais que tornam a sentença verdadeira.

Considere a seguinte inequação elementar do primeiro grau:

$ 4x - 18 < 2 $

Resolva a inequação no conjunto dos números Racionais ($\mathbb{Q}$) e indique a representação gráfica correta do seu conjunto solução na reta numérica real.
\end{enunciadoLiteral}

\ifgabarito
    % --- CAIXA LARANJA (Resolução e Tutor) ---
    \begin{tcolorbox}[colback=CorTitUm!5, colframe=CorTitUm, boxrule=1pt, arc=4pt, left=6pt, right=6pt, top=6pt, bottom=6pt]
    \small\textbf{\textcolor{CorTitUm}{Roteiro do Professor e Tutor IA:}}\par\vspace{0.1cm}
    \color{CinzaEscuro}
    \textbf{1. Ancoragem (Abordagem Socrática):} "O que o símbolo '$<$' significa? Se tivéssemos um sinal de igual, o que faríamos com o $-18$? Será que o princípio da balança ainda funciona para desigualdades se somarmos números positivos?"\par\vspace{0.1cm}
    
    \textbf{2. Micro-passos da Resolução:}\\
    \textit{Passo 1:} Princípio aditivo. Isolar o termo com incógnita adicionando 18 a ambos os membros $\rightarrow$ $4x < 2 + 18 \Rightarrow 4x < 20$.\\
    \textit{Passo 2:} Princípio multiplicativo. Dividir ambos os membros pelo número positivo 4 (o sentido da desigualdade permanece idêntico) $\rightarrow$ $x < \mfrac{20}{4} \Rightarrow x < 5$.\\
    \textit{Passo 3:} Conjunto Solução: $S = \{x \in \mathbb{Q} \mid x < 5\}$.\\
    \textit{Passo 4:} Construção na reta numérica: O ponto 5 deve receber bolinha aberta (não inclui o 5, pois é estritamente menor), com a semirreta destacada para a esquerda.\par\vspace{0.1cm}
    
    \textbf{3. Mapeamento de Armadilhas (Alerta Tutor):} O erro comum é o aluno preencher a bolinha no número 5 (confundindo menor estrito com menor ou igual), ou inverter o sentido da semirreta na reta numérica. Mostre que qualquer número menor que 5 (como 0 ou 4) satisfaz a sentença original ($4(0) - 18 = -18 < 2$).
    
    % RETA NUMÉRICA TIKZ DA RESOLUÇÃO
    \begin{center}
    \begin{adjustbox}{trim=0cm 0cm 0cm 0cm, clip, max width=0.85\linewidth}
    \begin{tikzpicture}
    % --- APLICANDO DROP SHADOW E CORES ---
    \definecolor{eixoReta}{HTML}{2C3E50}
    \definecolor{solucaoCor}{HTML}{E74C3C}
    
    % LAYER 1: Reta Numérica Base
    \draw[->, thick, eixoReta] (1,0) -- (8,0) node[right, font=\footnotesize\bfseries] {$\mathbb{R}$};
    
    % LAYER 2: Graduação
    \foreach \n in {2,3,4,5,6,7} {
        \draw[thick, eixoReta] (\n, 0.1) -- (\n, -0.1) node[below, font=\footnotesize] {\n};
    }
    
    % LAYER 3: Semirreta Solução (x < 5)
    \draw[ultra thick, solucaoCor] (1.2, 0) -- (5, 0);
    \fill[white] (5,0) circle (0.1);
    \draw[thick, solucaoCor] (5,0) circle (0.1);
    \node[fill=white, inner sep=1pt, text=solucaoCor, font=\bfseries\scriptsize] at (3, 0.4) {Região Solução: x < 5};
    \end{tikzpicture}
    \end{adjustbox}
    \end{center}
    \end{tcolorbox}
    \vspace{0.15cm}
    
    % --- CAIXA VERDE (Porto Seguro Didático do Professor) ---
    \begin{boxexplicacao}
    \textbf{Porto Seguro Didático (O "Pulo do Gato"):} Uma equação é um "ponto fixo" na reta. Uma inequação é um "mar aberto" de possibilidades. A bolinha aberta significa: "posso chegar perto de você (4,999...), mas não posso encostar no 5!". Se tivesse a barra embaixo ($\le$), a bolinha seria pintada de preto (fechada), permitindo pisar no número exato.
    \end{boxexplicacao}
\fi

\par\vspace{\espacoQuestao}


\noindent\textcolor{CorLinha}{\rule{\linewidth}{0.5pt}} 
\par\vspace{\espacoPreQuestao}
\subsubsection*{22.} 
\begin{enunciadoLiteral}
A propriedade mais crítica no estudo das inequações do 1º grau envolve a multiplicação ou divisão de seus membros por um número negativo.

Resolva a inequação a seguir, explicitando todo o passo a passo da regra de inversão de sinal da desigualdade:

$ 5 - 3x \ge 26 $
\end{enunciadoLiteral}

\ifgabarito
    % --- CAIXA LARANJA (Resolução e Tutor) ---
    \begin{tcolorbox}[colback=CorTitUm!5, colframe=CorTitUm, boxrule=1pt, arc=4pt, left=6pt, right=6pt, top=6pt, bottom=6pt]
    \small\textbf{\textcolor{CorTitUm}{Roteiro do Professor e Tutor IA:}}\par\vspace{0.1cm}
    \color{CinzaEscuro}
    \textbf{1. Ancoragem (Abordagem Socrática):} "Se $2 < 5$, o que acontece se multiplicarmos os dois números por $-1$? $-2$ é maior ou menor que $-5$? Na reta numérica, quem fica mais à direita?"\par\vspace{0.1cm}
    
    \textbf{2. Micro-passos da Resolução:}\\
    \textit{Passo 1:} Princípio aditivo. Subtrair 5 de ambos os membros $\rightarrow$ $-3x \ge 26 - 5 \Rightarrow -3x \ge 21$.\\
    \textit{Passo 2:} O Ponto Crítico da Inequação. Multiplicar toda a sentença por $-1$ (ou dividir diretamente por $-3$). Ao fazer isso, o sinal de desigualdade DEVE obrigatoriamente inverter de sentido $\rightarrow$ $3x \le -21$.\\
    \textit{Passo 3:} Princípio multiplicativo $\rightarrow$ $x \le \mfrac{-21}{3} \Rightarrow x \le -7$.\\
    \textit{Passo 4:} Conjunto Solução: $S = \{x \in \mathbb{Q} \mid x \le -7\}$.\par\vspace{0.1cm}
    
    \textbf{3. Mapeamento de Armadilhas (Alerta Tutor):} O erro número 1 em todo o Ensino Fundamental é o aluno manter o sinal inalterado ao passar o $-3$ dividindo, escrevendo $x \ge -7$. Faça o aluno testar o valor $x = 0$ (que é maior que $-7$): $5 - 3(0) = 5$, e 5 NÃO é maior ou igual a 26! Isso comprova a falha na hora.
    \end{tcolorbox}
    \vspace{0.15cm}
    
    % --- CAIXA VERDE (Porto Seguro Didático do Professor) ---
    \begin{boxexplicacao}
    \textbf{Porto Seguro Didático (O "Pulo do Gato"):} Pense na reta numérica como um espelho. O número 10 é maior que 2. Mas quando você olha no espelho dos negativos, o $-10$ está congelando muito mais longe do zero à esquerda, ou seja, $-10$ é menor que $-2$. Por isso, multiplicar por negativo inverte a hierarquia do universo numérico; a ponta da flecha tem que girar na mesma hora!
    \end{boxexplicacao}
\fi

\par\vspace{\espacoQuestao}


\noindent\textcolor{CorLinha}{\rule{\linewidth}{0.5pt}}
\par\vspace{\espacoPreQuestao}
\subsubsection*{23.} 
\begin{enunciadoLiteral}
Um estudante estava resolvendo a seguinte inequação no caderno:

$ 7 - 2(x + 4) > 13 $

Ele apresentou a seguinte sequência de linhas:

\textit{Linha 1:} $ 7 - 2x - 8 > 13 $ \\
\textit{Linha 2:} $ -2x - 1 > 13 $ \\
\textit{Linha 3:} $ -2x > 14 $ \\
\textit{Linha 4:} $ x > -7 $

O professor assinalou a Linha 4 com caneta vermelha, apontando um erro conceitual grave que comprometeu todo o resultado final.

\begin{boxresponda}
\textbf{Responda:} Identifique o erro cometido pelo estudante na transição da Linha 3 para a Linha 4. Explique textualmente a regra matemática violada e apresente o conjunto solução rigorosamente correto.
\end{boxresponda}
\end{enunciadoLiteral}

\ifgabarito
    % --- CAIXA LARANJA (Resolução e Tutor) ---
    \begin{tcolorbox}[colback=CorTitUm!5, colframe=CorTitUm, boxrule=1pt, arc=4pt, left=6pt, right=6pt, top=6pt, bottom=6pt]
    \small\textbf{\textcolor{CorTitUm}{Roteiro do Professor e Tutor IA:}}\par\vspace{0.1cm}
    \color{CinzaEscuro}
    \textbf{1. Ancoragem (Abordagem Socrática):} "O estudante fez tudo certinho até a Linha 3. Mas quando ele isolou o 'x' dividindo por $-2$, o que aconteceu com o bico da desigualdade?"\par\vspace{0.1cm}
    
    \textbf{2. Micro-passos da Resolução:}\\
    \textit{Passo 1:} Diagnóstico do erro. O estudante calculou a divisão corretamente ($14 / -2 = -7$), mas esqueceu de inverter o operador de desigualdade de '$>$' (maior que) para '$<$' (menor que).\\
    \textit{Passo 2:} Explicação conceitual. Toda vez que dividimos ou multiplicamos ambos os membros de uma desigualdade por uma grandeza estritamente negativa, a ordem da relação de grandeza se inverte.\\
    \textit{Passo 3:} Resolução correta a partir da Linha 3 $\rightarrow$ $-2x > 14 \Rightarrow x < \mfrac{14}{-2} \Rightarrow x < -7$.\\
    \textit{Passo 4:} Conjunto solução correto: $S = \{x \in \mathbb{Q} \mid x < -7\}$.\par\vspace{0.1cm}
    
    \textbf{3. Mapeamento de Armadilhas (Alerta Tutor):} Não permita que o aluno responda apenas "ele errou o sinal". Ele precisa dizer com clareza: "ele esqueceu de inverter o sentido da desigualdade". Precisão cirúrgica de vocabulário é o objetivo pedagógico aqui.
    \end{tcolorbox}
    \vspace{0.15cm}
    
    % --- CAIXA VERDE (Nota Pedagógica) ---
    \begin{boxexplicacao}
    \textbf{Justificativa Curricular (Raio-X da Questão):} Fixa a metacognição sobre a inversão do sinal em inequações, distinguindo operacionalmente o comportamento de equações (onde o sinal de igual nunca muda) e inequações.
    \end{boxexplicacao}
\fi

\par\vspace{\espacoQuestao}


\noindent\textcolor{CorLinha}{\rule{\linewidth}{0.5pt}} 
\par\vspace{\espacoPreQuestao}
\subsubsection*{24.} 
\begin{enunciadoLiteral}
A figura a seguir exibe a representação gráfica de um conjunto solução sobre a reta dos números inteiros ($\mathbb{Z}$).

\begin{center}
% ==========================================
% CONTROLE INDIVIDUAL DESTA IMAGEM:
% 1. CORTAR BORDAS: trim=0cm 0cm 0cm 0cm
% 2. TAMANHO: \larguraTikz por um valor
% ==========================================
\begin{adjustbox}{trim=0cm 0cm 0cm 0cm, clip, max width=0.85\linewidth}
\begin{tikzpicture}
% --- APLICANDO DROP SHADOW E CORES ---
\definecolor{eixoReta}{HTML}{2C3E50}
\definecolor{solucaoCor}{HTML}{27AE60}

% LAYER 1: Reta Numérica Base
\draw[->, thick, eixoReta] (-3,0) -- (5,0) node[right, font=\footnotesize\bfseries] {$\mathbb{Z}$};

% LAYER 2: Graduação de inteiros
\foreach \n in {-2,-1,0,1,2,3,4} {
    \draw[thick, eixoReta] (\n, 0.1) -- (\n, -0.1) node[below, font=\footnotesize] {\n};
}

% LAYER 3: Região destacada (x >= 1)
\draw[ultra thick, solucaoCor] (1, 0) -- (4.8, 0);
\fill[solucaoCor, drop shadow={opacity=0.2, shadow xshift=1pt, shadow yshift=-1pt}] (1,0) circle (0.12);
\node[fill=white, inner sep=2pt, rounded corners=1pt, text=solucaoCor, font=\bfseries\footnotesize] at (2.8, 0.45) {Intervalo Solução};
\end{tikzpicture}
\end{adjustbox}
\end{center}

Assinale a alternativa que contém a inequação do 1º grau cuja resolução produz rigorosamente a representação visual destacada na figura. Em seguida, prove por que a alternativa A está errada.

\begin{enumerate}[label=\alph*), itemsep=\espacoAlt]
\item $ 2x + 5 \le 7 $
\item $ 3x - 4 \ge -1 $
\item $ -2x + 8 < 6 $
\item $ 4x - 10 \ge 6 $
\end{enumerate}
\end{enunciadoLiteral}

\ifgabarito
    % --- CAIXA LARANJA (Resolução e Tutor) ---
    \begin{tcolorbox}[colback=CorTitUm!5, colframe=CorTitUm, boxrule=1pt, arc=4pt, left=6pt, right=6pt, top=6pt, bottom=6pt]
    \small\textbf{\textcolor{CorTitUm}{Roteiro do Professor e Tutor IA:}}\par\vspace{0.1cm}
    \color{CinzaEscuro}
    \textbf{1. Ancoragem (Abordagem Socrática):} "Olhando para a reta numérica: o ponto marcado é o número 1. A bolinha está cheia ou oca? A região pintada vai para a direita (maiores) ou para a esquerda (menores)? Qual condição matemática isso gera?"\par\vspace{0.1cm}
    
    \textbf{2. Micro-passos da Resolução:}\\
    \textit{Passo 1:} Leitura gráfica (Teste Cego). A bolinha preta cheia no número 1 com a flecha apontando para a direita representa a condição algébrica: $x \ge 1$.\\
    \textit{Passo 2:} Testando as alternativas:\\
    - Letra A: $2x + 5 \le 7 \Rightarrow 2x \le 2 \Rightarrow x \le 1$ (Sinal invertido, aponta para a esquerda).\\
    - Letra B: $3x - 4 \ge -1 \Rightarrow 3x \ge 3 \Rightarrow x \ge 1$ (Bate perfeitamente com a figura!).\\
    - Letra C: $-2x + 8 < 6 \Rightarrow -2x < -2 \Rightarrow x > 1$ (Bolinha seria aberta, não cheia).\\
    - Letra D: $4x - 10 \ge 6 \Rightarrow 4x \ge 16 \Rightarrow x \ge 4$ (Ponto inicial em 4).\\
    A alternativa correta é a \textbf{letra B}.\\
    \textit{Passo 3:} Refutação da Letra A. O erro da letra A está no sentido da desigualdade ($\le$). Embora o número 1 seja o ponto limite, a solução $x \le 1$ exige que os valores sejam menores que 1 (à esquerda), contradizendo o gráfico que avança para a direita.\par\vspace{0.1cm}
    
    \textbf{3. Mapeamento de Armadilhas (Alerta Tutor):} O aluno que resolve rápido costuma apenas olhar o número "1" e marcar a letra A sem reparar no símbolo de maior/menor. Chame a atenção para a ponta da semirreta colorida.
    \end{tcolorbox}
    \vspace{0.15cm}
    
    % --- CAIXA VERDE (Porto Seguro Didático do Professor) ---
    \begin{boxexplicacao}
    \textbf{Porto Seguro Didático (O "Pulo do Gato"):} Traduzir desenho para conta: se a flecha vai para a DIREITA, o 'x' quer crescer, então usamos maior que ($>$ ou $\ge$). Se a bolinha tem recheio de chocolate (pintada), tem tracinho de igual embaixo ($\ge$). Se for oca, é estritamente maior ($>$).
    \end{boxexplicacao}
\fi

\par\vspace{\espacoQuestao}


\noindent\textcolor{CorLinha}{\rule{\linewidth}{0.5pt}} 
\par\vspace{\espacoPreQuestao}
\subsubsection*{25.} 
\begin{enunciadoLiteral}
As inequações fracionárias exigem a mesma disciplina de agrupamento e desobstrução de denominadores que as equações.

Resolva a inequação a seguir no conjunto dos números Racionais ($\mathbb{Q}$), apresentando o conjunto solução formal:

$ \mfrac{x - 3}{2} - \mfrac{x + 1}{4} \le 1 $
\end{enunciadoLiteral}

\ifgabarito
    % --- CAIXA LARANJA (Resolução e Tutor) ---
    \begin{tcolorbox}[colback=CorTitUm!5, colframe=CorTitUm, boxrule=1pt, arc=4pt, left=6pt, right=6pt, top=6pt, bottom=6pt]
    \small\textbf{\textcolor{CorTitUm}{Roteiro do Professor e Tutor IA:}}\par\vspace{0.1cm}
    \color{CinzaEscuro}
    \textbf{1. Ancoragem (Abordagem Socrática):} "O que faremos com os denominadores 2 e 4? O MMC é 4. Como o número 4 é positivo, o sinal de desigualdade vai mudar de lado quando multiplicarmos a inequação inteira por ele?"\par\vspace{0.1cm}
    
    \textbf{2. Micro-passos da Resolução:}\\
    \textit{Passo 1:} O MMC entre 2, 4 e 1 é 4. Multiplicar todos os membros por 4 (sendo 4 positivo, a desigualdade se mantém inalterada) $\rightarrow$ $4 \cdot \left(\mfrac{x - 3}{2}\right) - 4 \cdot \left(\mfrac{x + 1}{4}\right) \le 4 \cdot (1)$.\\
    \textit{Passo 2:} Simplificar os coeficientes mantendo parênteses protetores nos numeradores $\rightarrow$ $2(x - 3) - 1(x + 1) \le 4$.\\
    \textit{Passo 3:} Distributiva cuidadosa com o sinal negativo do segundo bloco $\rightarrow$ $2x - 6 - x - 1 \le 4$.\\
    \textit{Passo 4:} Redução dos termos semelhantes $\rightarrow$ $x - 7 \le 4$.\\
    \textit{Passo 5:} Princípio aditivo final $\rightarrow$ $x \le 4 + 7 \Rightarrow x \le 11$.\\
    \textit{Passo 6:} Conjunto Solução: $S = \{x \in \mathbb{Q} \mid x \le 11\}$.\par\vspace{0.1cm}
    
    \textbf{3. Mapeamento de Armadilhas (Alerta Tutor):} Cuidado extremo com a fração $-\mfrac{x+1}{4}$. Quase todos os estudantes esquecem que o sinal de menos na frente da fração age como um parêntese e vira $-x - 1$, e não $-x + 1$. Aponte o lápis para o traço da fração.
    \end{tcolorbox}
    \vspace{0.15cm}
    
    % --- CAIXA VERDE (Porto Seguro Didático do Professor) ---
    \begin{boxexplicacao}
    \textbf{Porto Seguro Didático (O "Pulo do Gato"):} Recapitulando a regra de ouro das frações: o traço de uma fração funciona exatamente como uma corda amarrando todo mundo que está lá em cima. Se tiver um sinal negativo na frente dessa corda, quando você cortar o denominador pelo MMC, use parênteses protetores, senão o sinal de menos "engole" a sua nota!
    \end{boxexplicacao}
\fi

\par\vspace{\espacoQuestao}


\noindent\textcolor{CorLinha}{\rule{\linewidth}{0.5pt}} 
\par\vspace{\espacoPreQuestao}
\subsubsection*{26.} 
\begin{enunciadoLiteral}
Um elevador de carga em uma obra da construção civil possui limite rigoroso de segurança para evitar acidentes mecânicos no cabo de aço. O operador da obra sabe que o peso máximo que a cabine suporta sem entrar em sobrecarga é de \textbf{650 kg}.

O operador do elevador, que pesa exatamente \textbf{80 kg}, precisa transportar blocos de cimento estrutural. O diagrama técnico a seguir resume as cargas presentes na cabine.

\begin{center}
% ==========================================
% CONTROLE INDIVIDUAL DESTA IMAGEM:
% 1. CORTAR BORDAS: trim=0cm 0cm 0cm 0cm
% 2. TAMANHO: \larguraTikz por um valor
% ==========================================
\begin{adjustbox}{trim=0cm 0cm 0cm 0cm, clip, max width=0.75\linewidth}
\begin{tikzpicture}
% --- APLICANDO DROP SHADOW E CORES ---
\definecolor{cabineCor}{HTML}{34495E}
\definecolor{blocoCimento}{HTML}{7F8C8D}
\definecolor{operadorCor}{HTML}{E67E22}
\definecolor{avisoVermelho}{HTML}{C0392B}

% LAYER 1: Cabine do Elevador
\fill[gray!15, drop shadow={opacity=0.3, shadow xshift=2pt, shadow yshift=-2pt}, rounded corners=3pt] (0,0) rectangle (6,5);
\draw[cabineCor, ultra thick, rounded corners=3pt] (0,0) rectangle (6,5);

% Cabo de aço
\draw[line width=2pt, black] (3,5) -- (3,6.2);

% LAYER 2: Placa de Aviso no topo da cabine
\fill[avisoVermelho, rounded corners=2pt] (1.2, 4.2) rectangle (4.8, 4.7);
\node[text=white, font=\bfseries\footnotesize] at (3, 4.45) {CAPACIDADE MÁXIMA: 650 kg};

% LAYER 3: Operador (Esquerda)
\fill[operadorCor, rounded corners=2pt] (0.8, 0.5) rectangle (2.0, 3.2);
\node[text=white, font=\bfseries\small, align=center] at (1.4, 1.85) {Operador\\80 kg};

% LAYER 4: Pilha de Blocos de Cimento (Direita)
\fill[blocoCimento, drop shadow={opacity=0.2, shadow xshift=1pt, shadow yshift=-1pt}, rounded corners=1pt] (2.8, 0.5) rectangle (5.2, 1.2);
\fill[blocoCimento, drop shadow={opacity=0.2, shadow xshift=1pt, shadow yshift=-1pt}, rounded corners=1pt] (2.8, 1.3) rectangle (5.2, 2.0);
\fill[blocoCimento, drop shadow={opacity=0.2, shadow xshift=1pt, shadow yshift=-1pt}, rounded corners=1pt] (2.8, 2.1) rectangle (5.2, 2.8);
\node[fill=white, inner sep=2pt, rounded corners=1pt, font=\bfseries\footnotesize] at (4.0, 1.65) {N blocos de 25 kg cada};

\end{tikzpicture}
\end{adjustbox}
\end{center}

Monte a inequação do 1º grau que modela a segurança do elevador e descubra qual é o número máximo inteiro de blocos de cimento (\textbf{N}) que podem ser transportados em uma única viagem.
\end{enunciadoLiteral}

\ifgabarito
    % --- CAIXA LARANJA (Resolução e Tutor) ---
    \begin{tcolorbox}[colback=CorTitUm!5, colframe=CorTitUm, boxrule=1pt, arc=4pt, left=6pt, right=6pt, top=6pt, bottom=6pt]
    \small\textbf{\textcolor{CorTitUm}{Roteiro do Professor e Tutor IA:}}\par\vspace{0.1cm}
    \color{CinzaEscuro}
    \textbf{1. Ancoragem (Abordagem Socrática):} "O elevador aguenta 650 kg. Ele pode carregar 650 kg exatos? Sim. Pode carregar mais? Não. Qual símbolo matemático usamos para 'no máximo'? Como somamos o operador com os blocos?"\par\vspace{0.1cm}
    
    \textbf{2. Micro-passos da Resolução:}\\
    \textit{Passo 1:} Modelagem da carga total $\rightarrow$ Peso do Operador ($80$) $+$ Peso dos Blocos ($25N$).\\
    \textit{Passo 2:} A restrição de segurança "no máximo" traduz-se pelo símbolo de menor ou igual ($\le$) $\rightarrow$ $80 + 25N \le 650$.\\
    \textit{Passo 3:} Princípio aditivo $\rightarrow$ $25N \le 650 - 80 \Rightarrow 25N \le 570$.\\
    \textit{Passo 4:} Princípio multiplicativo $\rightarrow$ $N \le \mfrac{570}{25} \Rightarrow N \le 22,8$.\\
    \textit{Passo 5:} Análise contextual (Interpretação da Solução). Como não é possível carregar frações de bloco de cimento (o bloco precisa ser inteiro) e $N$ deve ser menor ou igual a 22,8, o maior número inteiro possível é $N = 22$ blocos.\par\vspace{0.1cm}
    
    \textbf{3. Mapeamento de Armadilhas (Alerta Tutor):} O erro clássico de arredondamento: o aluno olha 22,8 e pensa "vou arredondar para cima, 23!". O Tutor deve alertar: se colocar 23 blocos, a conta dá $80 + 25(23) = 655 \text{ kg}$, o que estoura a capacidade e rompe o elevador! Em inequações de teto máximo, arredonda-se para baixo para manter a margem de segurança.
    \end{tcolorbox}
    \vspace{0.15cm}
    
    % --- CAIXA VERDE (Porto Seguro Didático do Professor) ---
    \begin{boxexplicacao}
    \textbf{Porto Seguro Didático (O "Pulo do Gato"):} Cuidado com a vida real! Na matemática de papel, 22,8 fica perto de 23. Mas se você está enchendo a mala do avião ou o elevador de uma obra, ultrapassar o teto derruba o avião. Quando a inequação impõe um limite máximo ($\le$), qualquer número decimal precisa ser truncado para o inteiro imediatamente anterior.
    \end{boxexplicacao}
\fi

\par\vspace{\espacoQuestao}


\noindent\textcolor{CorLinha}{\rule{\linewidth}{0.5pt}} 
\par\vspace{\espacoPreQuestao}
\subsubsection*{27.} 
\begin{enunciadoLiteral}
(ENEM - Adaptada) Uma jovem autônoma decidiu vender bolos caseiros gourmet. O investimento inicial em formas e batedeira foi de R\$ 360,00 (custo fixo). Para produzir cada bolo, ela gasta exatamente R\$ 12,00 em ingredientes e embalagens. Ela fixou o preço de venda de cada bolo em R\$ 20,00.

Para que a doceira comece a obter **lucro real** (ou seja, o faturamento total supere rigorosamente todos os custos de produção somados ao investimento inicial), a quantidade mínima de bolos (\textbf{b}) que ela deve produzir e vender é dada por:

\begin{enumerate}[label=\alph*), itemsep=\espacoAlt]
\item $ b \ge 45 $
\item $ b \ge 46 $
\item $ b \le 45 $
\item $ b \ge 50 $
\end{enumerate}
\end{enunciadoLiteral}

\ifgabarito
    % --- CAIXA LARANJA (Resolução e Tutor) ---
    \begin{tcolorbox}[colback=CorTitUm!5, colframe=CorTitUm, boxrule=1pt, arc=4pt, left=6pt, right=6pt, top=6pt, bottom=6pt]
    \small\textbf{\textcolor{CorTitUm}{Roteiro do Professor e Tutor IA:}}\par\vspace{0.1cm}
    \color{CinzaEscuro}
    \textbf{1. Ancoragem (Abordagem Socrática):} "O que significa dar lucro em uma empresa? Faturar a mesma quantia que gastou é lucro ou apenas 'empate' (zero a zero)? Como escrevemos 'faturamento maior que o custo'?"\par\vspace{0.1cm}
    
    \textbf{2. Micro-passos da Resolução:}\\
    \textit{Passo 1:} Modelagem da Receita $\rightarrow$ $R = 20b$.\\
    \textit{Passo 2:} Modelagem do Custo Total $\rightarrow$ $C = 360 + 12b$.\\
    \textit{Passo 3:} Condição de lucro real: Receita maior que Custo $\rightarrow$ $20b > 360 + 12b$.\\
    \textit{Passo 4:} Resolução algébrica $\rightarrow$ $20b - 12b > 360 \Rightarrow 8b > 360$.\\
    \textit{Passo 5:} Divisão $\rightarrow$ $b > \mfrac{360}{8} \Rightarrow b > 45$.\\
    \textit{Passo 6:} Interpretação final. Se $b > 45$, produzir 45 bolos apenas empata o investimento ($20 \cdot 45 = 900$; $360 + 12 \cdot 45 = 900$). Para ter lucro real ($>0$), o primeiro número inteiro possível é 46 bolos. Logo, $b \ge 46$. Alternativa correta: \textbf{Letra B}.\par\vspace{0.1cm}
    
    \textbf{3. Mapeamento de Armadilhas (Alerta Tutor):} A alternativa A ($b \ge 45$) é o distrator fatal do ENEM. O aluno faz $360/8 = 45$ e marca a letra A no automático. O Tutor deve perguntar: "Se ela vender 45 bolos, quanto dinheiro sobra no bolso dela?". O aluno verá que dá zero e corrigirá para 46.
    \end{tcolorbox}
    \vspace{0.15cm}
    
    % --- CAIXA VERDE (Porto Seguro Didático do Professor) ---
    \begin{boxexplicacao}
    \textbf{Porto Seguro Didático (O "Pulo do Gato"):} "Empatar" não é "lucrar". O zero a zero apenas paga as contas. Em problemas econômicos de vestibular, a palavra 'lucro' exige o sinal de estritamente maior ($>$). Por isso, se a conta der fração ou número exato de equilíbrio, você dá um passo inteiro para a frente para sentir o primeiro centavo de lucro!
    \end{boxexplicacao}
\fi

\par\vspace{\espacoQuestao}


\noindent\textcolor{CorLinha}{\rule{\linewidth}{0.5pt}} 
\par\vspace{\espacoPreQuestao}
\subsubsection*{28.} 
\begin{enunciadoLiteral}
Uma operadora de telefonia disponibiliza um plano de dados pós-pago no valor fixo de R\$ 40,00 por mês, que dá direito a uma franquia base de internet. Se o cliente ultrapassar a franquia, é cobrada uma taxa de R\$ 2,50 por cada gigabyte (\textbf{g}) adicional consumido.

Um usuário planejou seu orçamento mensal e determinou que não pode gastar, sob hipótese alguma, mais do que R\$ 65,00 na sua fatura de celular neste mês.

Apresente a inequação que retrata a restrição financeira do usuário, determine a quantidade máxima de gigabytes adicionais que ele poderá consumir e construa a representação dessa solução na reta numérica real.
\end{enunciadoLiteral}

\ifgabarito
    % --- CAIXA LARANJA (Resolução e Tutor) ---
    \begin{tcolorbox}[colback=CorTitUm!5, colframe=CorTitUm, boxrule=1pt, arc=4pt, left=6pt, right=6pt, top=6pt, bottom=6pt]
    \small\textbf{\textcolor{CorTitUm}{Roteiro do Professor e Tutor IA:}}\par\vspace{0.1cm}
    \color{CinzaEscuro}
    \textbf{1. Ancoragem (Abordagem Socrática):} "O que significa 'não pode gastar mais do que 65 reais'? Ele pode gastar exatamente 65? Pode gastar menos? Que operador matemático representa essa frase?"\par\vspace{0.1cm}
    
    \textbf{2. Micro-passos da Resolução:}\\
    \textit{Passo 1:} Modelagem da despesa: Fixo ($40$) $+$ Adicionais ($2,50g$).\\
    \textit{Passo 2:} Restrição de gasto: despesa deve ser menor ou igual a 65 $\rightarrow$ $40 + 2,50g \le 65$.\\
    \textit{Passo 3:} Princípio aditivo $\rightarrow$ $2,50g \le 65 - 40 \Rightarrow 2,50g \le 25$.\\
    \textit{Passo 4:} Princípio multiplicativo $\rightarrow$ $g \le \mfrac{25}{2,50} \Rightarrow g \le 10 \text{ GB}$.\\
    \textit{Passo 5:} Limite inferior contextual: como o consumo de dados não pode ser negativo, temos $0 \le g \le 10$.\par\vspace{0.1cm}
    
    \textbf{3. Mapeamento de Armadilhas (Alerta Tutor):} Alunos tendem a ignorar o limite inferior zero, esquecendo que no mundo real grandezas físicas (como consumo de internet) não assumem valores negativos.
    
    % RETA NUMÉRICA TIKZ DA RESOLUÇÃO
    \begin{center}
    \begin{adjustbox}{trim=0cm 0cm 0cm 0cm, clip, max width=0.85\linewidth}
    \begin{tikzpicture}
    % --- APLICANDO DROP SHADOW E CORES ---
    \definecolor{eixoReta}{HTML}{2C3E50}
    \definecolor{solucaoCor}{HTML}{2980B9}
    
    % LAYER 1: Reta Numérica Base
    \draw[->, thick, eixoReta] (-1,0) -- (12,0) node[right, font=\footnotesize\bfseries] {$\mathbb{R}$};
    
    % LAYER 2: Graduação
    \foreach \n in {0,2,4,6,8,10} {
        \draw[thick, eixoReta] (\n, 0.1) -- (\n, -0.1) node[below, font=\footnotesize] {\n};
    }
    
    % LAYER 3: Intervalo Fechado [0, 10]
    \draw[ultra thick, solucaoCor] (0, 0) -- (10, 0);
    \fill[solucaoCor] (0,0) circle (0.12);
    \fill[solucaoCor] (10,0) circle (0.12);
    \node[fill=white, inner sep=1pt, text=solucaoCor, font=\bfseries\scriptsize] at (5, 0.4) {Limite de Uso: até 10 GB};
    \end{tikzpicture}
    \end{adjustbox}
    \end{center}
    \end{tcolorbox}
    \vspace{0.15cm}
    
    % --- CAIXA VERDE (Porto Seguro Didático do Professor) ---
    \begin{boxexplicacao}
    \textbf{Porto Seguro Didático (O "Pulo do Gato"):} "Não posso gastar mais que X" é o lema de quem tem juízo financeiro. Toda vez que você vir as expressões "no máximo", "até" ou "não superior a", use sem medo o sinal de $\le$ (menor ou igual). Ele coloca uma cerca intransponível no valor mais alto permitido.
    \end{boxexplicacao}
\fi

\par\vspace{\espacoQuestao}


\noindent\textcolor{CorLinha}{\rule{\linewidth}{0.5pt}} 
\par\vspace{\espacoPreQuestao}
\subsubsection*{29.} 
\begin{enunciadoLiteral}
Em geometria plana euclidiana, uma das propriedades fundamentais para que um triângulo exista fisicamente é a **desigualdade triangular**: a medida de qualquer lado deve ser estritamente menor do que a soma das medidas dos outros dois lados.

Considere um triângulo cujos lados medem \textbf{5 cm}, \textbf{9 cm} e o terceiro lado mede uma quantidade algébrica expressa por \textbf{2x + 1 cm}.

Sabendo que o lado \textbf{2x + 1} deve ser menor do que a soma dos outros dois, resolva a inequação correspondente e descubra o valor limite superior da constante \textbf{x}.
\end{enunciadoLiteral}

\ifgabarito
    % --- CAIXA LARANJA (Resolução e Tutor) ---
    \begin{tcolorbox}[colback=CorTitUm!5, colframe=CorTitUm, boxrule=1pt, arc=4pt, left=6pt, right=6pt, top=6pt, bottom=6pt]
    \small\textbf{\textcolor{CorTitUm}{Roteiro do Professor e Tutor IA:}}\par\vspace{0.1cm}
    \color{CinzaEscuro}
    \textbf{1. Ancoragem (Abordagem Socrática):} "O que o teorema diz? Um lado sozinho precisa ser MENOR que os outros dois somados. Como somamos 5 com 9? E como colocamos a expressão $2x+1$ menor do que essa soma?"\par\vspace{0.1cm}
    
    \textbf{2. Micro-passos da Resolução:}\\
    \textit{Passo 1:} Aplicação direta da regra geométrica $\rightarrow$ Lado $< \text{Lado}_1 + \text{Lado}_2$.\\
    \textit{Passo 2:} Montagem da Inequação $\rightarrow$ $2x + 1 < 5 + 9$.\\
    \textit{Passo 3:} Redução $\rightarrow$ $2x + 1 < 14$.\\
    \textit{Passo 4:} Princípio aditivo $\rightarrow$ $2x < 14 - 1 \Rightarrow 2x < 13$.\\
    \textit{Passo 5:} Princípio multiplicativo $\rightarrow$ $x < \mfrac{13}{2} \Rightarrow x < 6,5$.\\
    \textit{Passo 6:} Resposta: O valor limite superior de \textbf{x} deve ser estritamente menor do que $6,5 \text{ cm}$.\par\vspace{0.1cm}
    
    \textbf{3. Mapeamento de Armadilhas (Alerta Tutor):} O aluno costuma usar o sinal de menor ou igual ($\le$) por engano. Lembre-o de que se o lado for exatamente igual à soma dos outros dois ($2x+1 = 14$), o triângulo "amassa" e vira apenas uma linha reta esticada no chão, deixando de ser um polígono!
    \end{tcolorbox}
    \vspace{0.15cm}
    
    % --- CAIXA VERDE (Porto Seguro Didático do Professor) ---
    \begin{boxexplicacao}
    \textbf{Porto Seguro Didático (O "Pulo do Gato"):} A geometria conversa o tempo todo com a álgebra. Imagine dois gravetos de 5 e 9 cm ligados por uma ponta. A maior abertura possível entre eles é quando eles quase formam uma reta de 14 cm. Se o terceiro graveto tiver 14 cm ou mais, ele não fecha a figura de jeito nenhum! Por isso, o lado tem que ser estritamente menor que a soma.
    \end{boxexplicacao}
\fi

\par\vspace{\espacoQuestao}


\noindent\textcolor{CorLinha}{\rule{\linewidth}{0.5pt}} 
\par\vspace{\espacoPreQuestao}
\subsubsection*{30.} 
\begin{enunciadoLiteral}
(Estilo UNICAMP) Um laboratório farmacêutico desenvolveu uma estufa inteligente para o cultivo de bactérias benéficas usadas na produção de iogurtes. Para que a colônia sobreviva, a temperatura interna \textbf{T} (em graus Celsius) deve atender simultaneamente a um sistema de duas restrições ambientais:

\begin{enumerate}[label=\Roman*., itemsep=\espacoAlt]
\item A temperatura deve satisfazer a inequação linear: $ 3(T - 4) \ge 18 $
\item A temperatura não pode atingir valores de superaquecimento, devendo respeitar: $ 2T + 5 < 45 $
\end{enumerate}

Resolva as duas inequações de forma independente e determine o intervalo exato de temperaturas seguras no qual a estufa pode operar. Em seguida, responda: a temperatura de \textbf{20°C} é segura para a colônia?
\end{enunciadoLiteral}

\ifgabarito
    % --- CAIXA LARANJA (Resolução e Tutor) ---
    \begin{tcolorbox}[colback=CorTitUm!5, colframe=CorTitUm, boxrule=1pt, arc=4pt, left=6pt, right=6pt, top=6pt, bottom=6pt]
    \small\textbf{\textcolor{CorTitUm}{Roteiro do Professor e Tutor IA:}}\par\vspace{0.1cm}
    \color{CinzaEscuro}
    \textbf{1. Ancoragem (Abordagem Socrática):} "Temos duas regras que precisam acontecer ao mesmo tempo. Primeiro descobrimos o chão (a temperatura mínima) e depois o teto (a máxima). Como resolvemos cada uma separadamente?"\par\vspace{0.1cm}
    
    \textbf{2. Micro-passos da Resolução:}\\
    \textit{Passo 1:} Resolução da Restrição I (Piso térmico):\\
    $3(T - 4) \ge 18 \Rightarrow 3T - 12 \ge 18 \Rightarrow 3T \ge 30 \Rightarrow T \ge 10^\circ\text{C}$.\\
    \textit{Passo 2:} Resolução da Restrição II (Teto térmico):\\
    $2T + 5 < 45 \Rightarrow 2T < 40 \Rightarrow T < 20^\circ\text{C}$.\\
    \textit{Passo 3:} Intersecção dos intervalos: A temperatura deve ser maior ou igual a 10 e estritamente menor que 20 $\rightarrow$ $10^\circ\text{C} \le T < 20^\circ\text{C}$.\\
    \textit{Passo 4:} Julgamento da pergunta final: A temperatura de \textbf{20°C} NÃO é segura, pois a restrição II exige $T < 20^\circ\text{C}$ (desigualdade estrita). Em 20°C ocorre superaquecimento e morte das bactérias.\par\vspace{0.1cm}
    
    \textbf{3. Mapeamento de Armadilhas (Alerta Tutor):} O aluno descuidado dirá que "20°C é seguro porque deu 20 na conta". Mostre o símbolo '$<$' e reforce que 20 não é menor do que 20.
    \end{tcolorbox}
    \vspace{0.15cm}
    
    % --- CAIXA VERDE (Porto Seguro Didático do Professor) ---
    \begin{boxexplicacao}
    \textbf{Porto Seguro Didático (O "Pulo do Gato"):} Quando temos duas condições juntas, estamos criando uma "faixa de sobrevivência" (intersecção). O 10°C entra na festa porque tem o traço de igual ($\ge$). Mas o 20°C fica barrado na porta porque a desigualdade é estrita ($<$). No vestibular, esse detalhe entre bolinha aberta e fechada vale a vaga!
    \end{boxexplicacao}
\fi

\par\vspace{\espacoQuestao}

\end{multicols*}
\end{document}